\documentclass[11pt]{article}

    \usepackage[breakable]{tcolorbox}
    \usepackage{parskip} % Stop auto-indenting (to mimic markdown behaviour)
    

    % Basic figure setup, for now with no caption control since it's done
    % automatically by Pandoc (which extracts ![](path) syntax from Markdown).
    \usepackage{graphicx}
    % Keep aspect ratio if custom image width or height is specified
    \setkeys{Gin}{keepaspectratio}
    % Maintain compatibility with old templates. Remove in nbconvert 6.0
    \let\Oldincludegraphics\includegraphics
    % Ensure that by default, figures have no caption (until we provide a
    % proper Figure object with a Caption API and a way to capture that
    % in the conversion process - todo).
    \usepackage{caption}
    \DeclareCaptionFormat{nocaption}{}
    \captionsetup{format=nocaption,aboveskip=0pt,belowskip=0pt}

    \usepackage{float}
    \floatplacement{figure}{H} % forces figures to be placed at the correct location
    \usepackage{xcolor} % Allow colors to be defined
    \usepackage{enumerate} % Needed for markdown enumerations to work
    \usepackage{geometry} % Used to adjust the document margins
    \usepackage{amsmath} % Equations
    \usepackage{amssymb} % Equations
    \usepackage{textcomp} % defines textquotesingle
    % Hack from http://tex.stackexchange.com/a/47451/13684:
    \AtBeginDocument{%
        \def\PYZsq{\textquotesingle}% Upright quotes in Pygmentized code
    }
    \usepackage{upquote} % Upright quotes for verbatim code
    \usepackage{eurosym} % defines \euro

    \usepackage{iftex}
    \ifPDFTeX
        \usepackage[T1]{fontenc}
        \IfFileExists{alphabeta.sty}{
              \usepackage{alphabeta}
          }{
              \usepackage[mathletters]{ucs}
              \usepackage[utf8x]{inputenc}
          }
    \else
        \usepackage{fontspec}
        \usepackage{unicode-math}
    \fi

    \usepackage{fancyvrb} % verbatim replacement that allows latex
    \usepackage{grffile} % extends the file name processing of package graphics
                         % to support a larger range
    \makeatletter % fix for old versions of grffile with XeLaTeX
    \@ifpackagelater{grffile}{2019/11/01}
    {
      % Do nothing on new versions
    }
    {
      \def\Gread@@xetex#1{%
        \IfFileExists{"\Gin@base".bb}%
        {\Gread@eps{\Gin@base.bb}}%
        {\Gread@@xetex@aux#1}%
      }
    }
    \makeatother
    \usepackage[Export]{adjustbox} % Used to constrain images to a maximum size
    \adjustboxset{max size={0.9\linewidth}{0.9\paperheight}}

    % The hyperref package gives us a pdf with properly built
    % internal navigation ('pdf bookmarks' for the table of contents,
    % internal cross-reference links, web links for URLs, etc.)
    \usepackage{hyperref}
    % The default LaTeX title has an obnoxious amount of whitespace. By default,
    % titling removes some of it. It also provides customization options.
    \usepackage{titling}
    \usepackage{longtable} % longtable support required by pandoc >1.10
    \usepackage{booktabs}  % table support for pandoc > 1.12.2
    \usepackage{array}     % table support for pandoc >= 2.11.3
    \usepackage{calc}      % table minipage width calculation for pandoc >= 2.11.1
    \usepackage[inline]{enumitem} % IRkernel/repr support (it uses the enumerate* environment)
    \usepackage[normalem]{ulem} % ulem is needed to support strikethroughs (\sout)
                                % normalem makes italics be italics, not underlines
    \usepackage{soul}      % strikethrough (\st) support for pandoc >= 3.0.0
    \usepackage{mathrsfs}
    

    
    % Colors for the hyperref package
    \definecolor{urlcolor}{rgb}{0,.145,.698}
    \definecolor{linkcolor}{rgb}{.71,0.21,0.01}
    \definecolor{citecolor}{rgb}{.12,.54,.11}

    % ANSI colors
    \definecolor{ansi-black}{HTML}{3E424D}
    \definecolor{ansi-black-intense}{HTML}{282C36}
    \definecolor{ansi-red}{HTML}{E75C58}
    \definecolor{ansi-red-intense}{HTML}{B22B31}
    \definecolor{ansi-green}{HTML}{00A250}
    \definecolor{ansi-green-intense}{HTML}{007427}
    \definecolor{ansi-yellow}{HTML}{DDB62B}
    \definecolor{ansi-yellow-intense}{HTML}{B27D12}
    \definecolor{ansi-blue}{HTML}{208FFB}
    \definecolor{ansi-blue-intense}{HTML}{0065CA}
    \definecolor{ansi-magenta}{HTML}{D160C4}
    \definecolor{ansi-magenta-intense}{HTML}{A03196}
    \definecolor{ansi-cyan}{HTML}{60C6C8}
    \definecolor{ansi-cyan-intense}{HTML}{258F8F}
    \definecolor{ansi-white}{HTML}{C5C1B4}
    \definecolor{ansi-white-intense}{HTML}{A1A6B2}
    \definecolor{ansi-default-inverse-fg}{HTML}{FFFFFF}
    \definecolor{ansi-default-inverse-bg}{HTML}{000000}

    % common color for the border for error outputs.
    \definecolor{outerrorbackground}{HTML}{FFDFDF}

    % commands and environments needed by pandoc snippets
    % extracted from the output of `pandoc -s`
    \providecommand{\tightlist}{%
      \setlength{\itemsep}{0pt}\setlength{\parskip}{0pt}}
    \DefineVerbatimEnvironment{Highlighting}{Verbatim}{commandchars=\\\{\}}
    % Add ',fontsize=\small' for more characters per line
    \newenvironment{Shaded}{}{}
    \newcommand{\KeywordTok}[1]{\textcolor[rgb]{0.00,0.44,0.13}{\textbf{{#1}}}}
    \newcommand{\DataTypeTok}[1]{\textcolor[rgb]{0.56,0.13,0.00}{{#1}}}
    \newcommand{\DecValTok}[1]{\textcolor[rgb]{0.25,0.63,0.44}{{#1}}}
    \newcommand{\BaseNTok}[1]{\textcolor[rgb]{0.25,0.63,0.44}{{#1}}}
    \newcommand{\FloatTok}[1]{\textcolor[rgb]{0.25,0.63,0.44}{{#1}}}
    \newcommand{\CharTok}[1]{\textcolor[rgb]{0.25,0.44,0.63}{{#1}}}
    \newcommand{\StringTok}[1]{\textcolor[rgb]{0.25,0.44,0.63}{{#1}}}
    \newcommand{\CommentTok}[1]{\textcolor[rgb]{0.38,0.63,0.69}{\textit{{#1}}}}
    \newcommand{\OtherTok}[1]{\textcolor[rgb]{0.00,0.44,0.13}{{#1}}}
    \newcommand{\AlertTok}[1]{\textcolor[rgb]{1.00,0.00,0.00}{\textbf{{#1}}}}
    \newcommand{\FunctionTok}[1]{\textcolor[rgb]{0.02,0.16,0.49}{{#1}}}
    \newcommand{\RegionMarkerTok}[1]{{#1}}
    \newcommand{\ErrorTok}[1]{\textcolor[rgb]{1.00,0.00,0.00}{\textbf{{#1}}}}
    \newcommand{\NormalTok}[1]{{#1}}

    % Additional commands for more recent versions of Pandoc
    \newcommand{\ConstantTok}[1]{\textcolor[rgb]{0.53,0.00,0.00}{{#1}}}
    \newcommand{\SpecialCharTok}[1]{\textcolor[rgb]{0.25,0.44,0.63}{{#1}}}
    \newcommand{\VerbatimStringTok}[1]{\textcolor[rgb]{0.25,0.44,0.63}{{#1}}}
    \newcommand{\SpecialStringTok}[1]{\textcolor[rgb]{0.73,0.40,0.53}{{#1}}}
    \newcommand{\ImportTok}[1]{{#1}}
    \newcommand{\DocumentationTok}[1]{\textcolor[rgb]{0.73,0.13,0.13}{\textit{{#1}}}}
    \newcommand{\AnnotationTok}[1]{\textcolor[rgb]{0.38,0.63,0.69}{\textbf{\textit{{#1}}}}}
    \newcommand{\CommentVarTok}[1]{\textcolor[rgb]{0.38,0.63,0.69}{\textbf{\textit{{#1}}}}}
    \newcommand{\VariableTok}[1]{\textcolor[rgb]{0.10,0.09,0.49}{{#1}}}
    \newcommand{\ControlFlowTok}[1]{\textcolor[rgb]{0.00,0.44,0.13}{\textbf{{#1}}}}
    \newcommand{\OperatorTok}[1]{\textcolor[rgb]{0.40,0.40,0.40}{{#1}}}
    \newcommand{\BuiltInTok}[1]{{#1}}
    \newcommand{\ExtensionTok}[1]{{#1}}
    \newcommand{\PreprocessorTok}[1]{\textcolor[rgb]{0.74,0.48,0.00}{{#1}}}
    \newcommand{\AttributeTok}[1]{\textcolor[rgb]{0.49,0.56,0.16}{{#1}}}
    \newcommand{\InformationTok}[1]{\textcolor[rgb]{0.38,0.63,0.69}{\textbf{\textit{{#1}}}}}
    \newcommand{\WarningTok}[1]{\textcolor[rgb]{0.38,0.63,0.69}{\textbf{\textit{{#1}}}}}
    \makeatletter
    \newsavebox\pandoc@box
    \newcommand*\pandocbounded[1]{%
      \sbox\pandoc@box{#1}%
      % scaling factors for width and height
      \Gscale@div\@tempa\textheight{\dimexpr\ht\pandoc@box+\dp\pandoc@box\relax}%
      \Gscale@div\@tempb\linewidth{\wd\pandoc@box}%
      % select the smaller of both
      \ifdim\@tempb\p@<\@tempa\p@
        \let\@tempa\@tempb
      \fi
      % scaling accordingly (\@tempa < 1)
      \ifdim\@tempa\p@<\p@
        \scalebox{\@tempa}{\usebox\pandoc@box}%
      % scaling not needed, use as it is
      \else
        \usebox{\pandoc@box}%
      \fi
    }
    \makeatother

    % Define a nice break command that doesn't care if a line doesn't already
    % exist.
    \def\br{\hspace*{\fill} \\* }
    % Math Jax compatibility definitions
    \def\gt{>}
    \def\lt{<}
    \let\Oldtex\TeX
    \let\Oldlatex\LaTeX
    \renewcommand{\TeX}{\textrm{\Oldtex}}
    \renewcommand{\LaTeX}{\textrm{\Oldlatex}}
    % Document parameters
    % Document title
    \title{Ch07-nonlin-lab}
    
    
    
    
    
    
    
% Pygments definitions
\makeatletter
\def\PY@reset{\let\PY@it=\relax \let\PY@bf=\relax%
    \let\PY@ul=\relax \let\PY@tc=\relax%
    \let\PY@bc=\relax \let\PY@ff=\relax}
\def\PY@tok#1{\csname PY@tok@#1\endcsname}
\def\PY@toks#1+{\ifx\relax#1\empty\else%
    \PY@tok{#1}\expandafter\PY@toks\fi}
\def\PY@do#1{\PY@bc{\PY@tc{\PY@ul{%
    \PY@it{\PY@bf{\PY@ff{#1}}}}}}}
\def\PY#1#2{\PY@reset\PY@toks#1+\relax+\PY@do{#2}}

\@namedef{PY@tok@w}{\def\PY@tc##1{\textcolor[rgb]{0.73,0.73,0.73}{##1}}}
\@namedef{PY@tok@c}{\let\PY@it=\textit\def\PY@tc##1{\textcolor[rgb]{0.24,0.48,0.48}{##1}}}
\@namedef{PY@tok@cp}{\def\PY@tc##1{\textcolor[rgb]{0.61,0.40,0.00}{##1}}}
\@namedef{PY@tok@k}{\let\PY@bf=\textbf\def\PY@tc##1{\textcolor[rgb]{0.00,0.50,0.00}{##1}}}
\@namedef{PY@tok@kp}{\def\PY@tc##1{\textcolor[rgb]{0.00,0.50,0.00}{##1}}}
\@namedef{PY@tok@kt}{\def\PY@tc##1{\textcolor[rgb]{0.69,0.00,0.25}{##1}}}
\@namedef{PY@tok@o}{\def\PY@tc##1{\textcolor[rgb]{0.40,0.40,0.40}{##1}}}
\@namedef{PY@tok@ow}{\let\PY@bf=\textbf\def\PY@tc##1{\textcolor[rgb]{0.67,0.13,1.00}{##1}}}
\@namedef{PY@tok@nb}{\def\PY@tc##1{\textcolor[rgb]{0.00,0.50,0.00}{##1}}}
\@namedef{PY@tok@nf}{\def\PY@tc##1{\textcolor[rgb]{0.00,0.00,1.00}{##1}}}
\@namedef{PY@tok@nc}{\let\PY@bf=\textbf\def\PY@tc##1{\textcolor[rgb]{0.00,0.00,1.00}{##1}}}
\@namedef{PY@tok@nn}{\let\PY@bf=\textbf\def\PY@tc##1{\textcolor[rgb]{0.00,0.00,1.00}{##1}}}
\@namedef{PY@tok@ne}{\let\PY@bf=\textbf\def\PY@tc##1{\textcolor[rgb]{0.80,0.25,0.22}{##1}}}
\@namedef{PY@tok@nv}{\def\PY@tc##1{\textcolor[rgb]{0.10,0.09,0.49}{##1}}}
\@namedef{PY@tok@no}{\def\PY@tc##1{\textcolor[rgb]{0.53,0.00,0.00}{##1}}}
\@namedef{PY@tok@nl}{\def\PY@tc##1{\textcolor[rgb]{0.46,0.46,0.00}{##1}}}
\@namedef{PY@tok@ni}{\let\PY@bf=\textbf\def\PY@tc##1{\textcolor[rgb]{0.44,0.44,0.44}{##1}}}
\@namedef{PY@tok@na}{\def\PY@tc##1{\textcolor[rgb]{0.41,0.47,0.13}{##1}}}
\@namedef{PY@tok@nt}{\let\PY@bf=\textbf\def\PY@tc##1{\textcolor[rgb]{0.00,0.50,0.00}{##1}}}
\@namedef{PY@tok@nd}{\def\PY@tc##1{\textcolor[rgb]{0.67,0.13,1.00}{##1}}}
\@namedef{PY@tok@s}{\def\PY@tc##1{\textcolor[rgb]{0.73,0.13,0.13}{##1}}}
\@namedef{PY@tok@sd}{\let\PY@it=\textit\def\PY@tc##1{\textcolor[rgb]{0.73,0.13,0.13}{##1}}}
\@namedef{PY@tok@si}{\let\PY@bf=\textbf\def\PY@tc##1{\textcolor[rgb]{0.64,0.35,0.47}{##1}}}
\@namedef{PY@tok@se}{\let\PY@bf=\textbf\def\PY@tc##1{\textcolor[rgb]{0.67,0.36,0.12}{##1}}}
\@namedef{PY@tok@sr}{\def\PY@tc##1{\textcolor[rgb]{0.64,0.35,0.47}{##1}}}
\@namedef{PY@tok@ss}{\def\PY@tc##1{\textcolor[rgb]{0.10,0.09,0.49}{##1}}}
\@namedef{PY@tok@sx}{\def\PY@tc##1{\textcolor[rgb]{0.00,0.50,0.00}{##1}}}
\@namedef{PY@tok@m}{\def\PY@tc##1{\textcolor[rgb]{0.40,0.40,0.40}{##1}}}
\@namedef{PY@tok@gh}{\let\PY@bf=\textbf\def\PY@tc##1{\textcolor[rgb]{0.00,0.00,0.50}{##1}}}
\@namedef{PY@tok@gu}{\let\PY@bf=\textbf\def\PY@tc##1{\textcolor[rgb]{0.50,0.00,0.50}{##1}}}
\@namedef{PY@tok@gd}{\def\PY@tc##1{\textcolor[rgb]{0.63,0.00,0.00}{##1}}}
\@namedef{PY@tok@gi}{\def\PY@tc##1{\textcolor[rgb]{0.00,0.52,0.00}{##1}}}
\@namedef{PY@tok@gr}{\def\PY@tc##1{\textcolor[rgb]{0.89,0.00,0.00}{##1}}}
\@namedef{PY@tok@ge}{\let\PY@it=\textit}
\@namedef{PY@tok@gs}{\let\PY@bf=\textbf}
\@namedef{PY@tok@ges}{\let\PY@bf=\textbf\let\PY@it=\textit}
\@namedef{PY@tok@gp}{\let\PY@bf=\textbf\def\PY@tc##1{\textcolor[rgb]{0.00,0.00,0.50}{##1}}}
\@namedef{PY@tok@go}{\def\PY@tc##1{\textcolor[rgb]{0.44,0.44,0.44}{##1}}}
\@namedef{PY@tok@gt}{\def\PY@tc##1{\textcolor[rgb]{0.00,0.27,0.87}{##1}}}
\@namedef{PY@tok@err}{\def\PY@bc##1{{\setlength{\fboxsep}{\string -\fboxrule}\fcolorbox[rgb]{1.00,0.00,0.00}{1,1,1}{\strut ##1}}}}
\@namedef{PY@tok@kc}{\let\PY@bf=\textbf\def\PY@tc##1{\textcolor[rgb]{0.00,0.50,0.00}{##1}}}
\@namedef{PY@tok@kd}{\let\PY@bf=\textbf\def\PY@tc##1{\textcolor[rgb]{0.00,0.50,0.00}{##1}}}
\@namedef{PY@tok@kn}{\let\PY@bf=\textbf\def\PY@tc##1{\textcolor[rgb]{0.00,0.50,0.00}{##1}}}
\@namedef{PY@tok@kr}{\let\PY@bf=\textbf\def\PY@tc##1{\textcolor[rgb]{0.00,0.50,0.00}{##1}}}
\@namedef{PY@tok@bp}{\def\PY@tc##1{\textcolor[rgb]{0.00,0.50,0.00}{##1}}}
\@namedef{PY@tok@fm}{\def\PY@tc##1{\textcolor[rgb]{0.00,0.00,1.00}{##1}}}
\@namedef{PY@tok@vc}{\def\PY@tc##1{\textcolor[rgb]{0.10,0.09,0.49}{##1}}}
\@namedef{PY@tok@vg}{\def\PY@tc##1{\textcolor[rgb]{0.10,0.09,0.49}{##1}}}
\@namedef{PY@tok@vi}{\def\PY@tc##1{\textcolor[rgb]{0.10,0.09,0.49}{##1}}}
\@namedef{PY@tok@vm}{\def\PY@tc##1{\textcolor[rgb]{0.10,0.09,0.49}{##1}}}
\@namedef{PY@tok@sa}{\def\PY@tc##1{\textcolor[rgb]{0.73,0.13,0.13}{##1}}}
\@namedef{PY@tok@sb}{\def\PY@tc##1{\textcolor[rgb]{0.73,0.13,0.13}{##1}}}
\@namedef{PY@tok@sc}{\def\PY@tc##1{\textcolor[rgb]{0.73,0.13,0.13}{##1}}}
\@namedef{PY@tok@dl}{\def\PY@tc##1{\textcolor[rgb]{0.73,0.13,0.13}{##1}}}
\@namedef{PY@tok@s2}{\def\PY@tc##1{\textcolor[rgb]{0.73,0.13,0.13}{##1}}}
\@namedef{PY@tok@sh}{\def\PY@tc##1{\textcolor[rgb]{0.73,0.13,0.13}{##1}}}
\@namedef{PY@tok@s1}{\def\PY@tc##1{\textcolor[rgb]{0.73,0.13,0.13}{##1}}}
\@namedef{PY@tok@mb}{\def\PY@tc##1{\textcolor[rgb]{0.40,0.40,0.40}{##1}}}
\@namedef{PY@tok@mf}{\def\PY@tc##1{\textcolor[rgb]{0.40,0.40,0.40}{##1}}}
\@namedef{PY@tok@mh}{\def\PY@tc##1{\textcolor[rgb]{0.40,0.40,0.40}{##1}}}
\@namedef{PY@tok@mi}{\def\PY@tc##1{\textcolor[rgb]{0.40,0.40,0.40}{##1}}}
\@namedef{PY@tok@il}{\def\PY@tc##1{\textcolor[rgb]{0.40,0.40,0.40}{##1}}}
\@namedef{PY@tok@mo}{\def\PY@tc##1{\textcolor[rgb]{0.40,0.40,0.40}{##1}}}
\@namedef{PY@tok@ch}{\let\PY@it=\textit\def\PY@tc##1{\textcolor[rgb]{0.24,0.48,0.48}{##1}}}
\@namedef{PY@tok@cm}{\let\PY@it=\textit\def\PY@tc##1{\textcolor[rgb]{0.24,0.48,0.48}{##1}}}
\@namedef{PY@tok@cpf}{\let\PY@it=\textit\def\PY@tc##1{\textcolor[rgb]{0.24,0.48,0.48}{##1}}}
\@namedef{PY@tok@c1}{\let\PY@it=\textit\def\PY@tc##1{\textcolor[rgb]{0.24,0.48,0.48}{##1}}}
\@namedef{PY@tok@cs}{\let\PY@it=\textit\def\PY@tc##1{\textcolor[rgb]{0.24,0.48,0.48}{##1}}}

\def\PYZbs{\char`\\}
\def\PYZus{\char`\_}
\def\PYZob{\char`\{}
\def\PYZcb{\char`\}}
\def\PYZca{\char`\^}
\def\PYZam{\char`\&}
\def\PYZlt{\char`\<}
\def\PYZgt{\char`\>}
\def\PYZsh{\char`\#}
\def\PYZpc{\char`\%}
\def\PYZdl{\char`\$}
\def\PYZhy{\char`\-}
\def\PYZsq{\char`\'}
\def\PYZdq{\char`\"}
\def\PYZti{\char`\~}
% for compatibility with earlier versions
\def\PYZat{@}
\def\PYZlb{[}
\def\PYZrb{]}
\makeatother


    % For linebreaks inside Verbatim environment from package fancyvrb.
    \makeatletter
        \newbox\Wrappedcontinuationbox
        \newbox\Wrappedvisiblespacebox
        \newcommand*\Wrappedvisiblespace {\textcolor{red}{\textvisiblespace}}
        \newcommand*\Wrappedcontinuationsymbol {\textcolor{red}{\llap{\tiny$\m@th\hookrightarrow$}}}
        \newcommand*\Wrappedcontinuationindent {3ex }
        \newcommand*\Wrappedafterbreak {\kern\Wrappedcontinuationindent\copy\Wrappedcontinuationbox}
        % Take advantage of the already applied Pygments mark-up to insert
        % potential linebreaks for TeX processing.
        %        {, <, #, %, $, ' and ": go to next line.
        %        _, }, ^, &, >, - and ~: stay at end of broken line.
        % Use of \textquotesingle for straight quote.
        \newcommand*\Wrappedbreaksatspecials {%
            \def\PYGZus{\discretionary{\char`\_}{\Wrappedafterbreak}{\char`\_}}%
            \def\PYGZob{\discretionary{}{\Wrappedafterbreak\char`\{}{\char`\{}}%
            \def\PYGZcb{\discretionary{\char`\}}{\Wrappedafterbreak}{\char`\}}}%
            \def\PYGZca{\discretionary{\char`\^}{\Wrappedafterbreak}{\char`\^}}%
            \def\PYGZam{\discretionary{\char`\&}{\Wrappedafterbreak}{\char`\&}}%
            \def\PYGZlt{\discretionary{}{\Wrappedafterbreak\char`\<}{\char`\<}}%
            \def\PYGZgt{\discretionary{\char`\>}{\Wrappedafterbreak}{\char`\>}}%
            \def\PYGZsh{\discretionary{}{\Wrappedafterbreak\char`\#}{\char`\#}}%
            \def\PYGZpc{\discretionary{}{\Wrappedafterbreak\char`\%}{\char`\%}}%
            \def\PYGZdl{\discretionary{}{\Wrappedafterbreak\char`\$}{\char`\$}}%
            \def\PYGZhy{\discretionary{\char`\-}{\Wrappedafterbreak}{\char`\-}}%
            \def\PYGZsq{\discretionary{}{\Wrappedafterbreak\textquotesingle}{\textquotesingle}}%
            \def\PYGZdq{\discretionary{}{\Wrappedafterbreak\char`\"}{\char`\"}}%
            \def\PYGZti{\discretionary{\char`\~}{\Wrappedafterbreak}{\char`\~}}%
        }
        % Some characters . , ; ? ! / are not pygmentized.
        % This macro makes them "active" and they will insert potential linebreaks
        \newcommand*\Wrappedbreaksatpunct {%
            \lccode`\~`\.\lowercase{\def~}{\discretionary{\hbox{\char`\.}}{\Wrappedafterbreak}{\hbox{\char`\.}}}%
            \lccode`\~`\,\lowercase{\def~}{\discretionary{\hbox{\char`\,}}{\Wrappedafterbreak}{\hbox{\char`\,}}}%
            \lccode`\~`\;\lowercase{\def~}{\discretionary{\hbox{\char`\;}}{\Wrappedafterbreak}{\hbox{\char`\;}}}%
            \lccode`\~`\:\lowercase{\def~}{\discretionary{\hbox{\char`\:}}{\Wrappedafterbreak}{\hbox{\char`\:}}}%
            \lccode`\~`\?\lowercase{\def~}{\discretionary{\hbox{\char`\?}}{\Wrappedafterbreak}{\hbox{\char`\?}}}%
            \lccode`\~`\!\lowercase{\def~}{\discretionary{\hbox{\char`\!}}{\Wrappedafterbreak}{\hbox{\char`\!}}}%
            \lccode`\~`\/\lowercase{\def~}{\discretionary{\hbox{\char`\/}}{\Wrappedafterbreak}{\hbox{\char`\/}}}%
            \catcode`\.\active
            \catcode`\,\active
            \catcode`\;\active
            \catcode`\:\active
            \catcode`\?\active
            \catcode`\!\active
            \catcode`\/\active
            \lccode`\~`\~
        }
    \makeatother

    \let\OriginalVerbatim=\Verbatim
    \makeatletter
    \renewcommand{\Verbatim}[1][1]{%
        %\parskip\z@skip
        \sbox\Wrappedcontinuationbox {\Wrappedcontinuationsymbol}%
        \sbox\Wrappedvisiblespacebox {\FV@SetupFont\Wrappedvisiblespace}%
        \def\FancyVerbFormatLine ##1{\hsize\linewidth
            \vtop{\raggedright\hyphenpenalty\z@\exhyphenpenalty\z@
                \doublehyphendemerits\z@\finalhyphendemerits\z@
                \strut ##1\strut}%
        }%
        % If the linebreak is at a space, the latter will be displayed as visible
        % space at end of first line, and a continuation symbol starts next line.
        % Stretch/shrink are however usually zero for typewriter font.
        \def\FV@Space {%
            \nobreak\hskip\z@ plus\fontdimen3\font minus\fontdimen4\font
            \discretionary{\copy\Wrappedvisiblespacebox}{\Wrappedafterbreak}
            {\kern\fontdimen2\font}%
        }%

        % Allow breaks at special characters using \PYG... macros.
        \Wrappedbreaksatspecials
        % Breaks at punctuation characters . , ; ? ! and / need catcode=\active
        \OriginalVerbatim[#1,codes*=\Wrappedbreaksatpunct]%
    }
    \makeatother

    % Exact colors from NB
    \definecolor{incolor}{HTML}{303F9F}
    \definecolor{outcolor}{HTML}{D84315}
    \definecolor{cellborder}{HTML}{CFCFCF}
    \definecolor{cellbackground}{HTML}{F7F7F7}

    % prompt
    \makeatletter
    \newcommand{\boxspacing}{\kern\kvtcb@left@rule\kern\kvtcb@boxsep}
    \makeatother
    \newcommand{\prompt}[4]{
        {\ttfamily\llap{{\color{#2}[#3]:\hspace{3pt}#4}}\vspace{-\baselineskip}}
    }
    

    
    % Prevent overflowing lines due to hard-to-break entities
    \sloppy
    % Setup hyperref package
    \hypersetup{
      breaklinks=true,  % so long urls are correctly broken across lines
      colorlinks=true,
      urlcolor=urlcolor,
      linkcolor=linkcolor,
      citecolor=citecolor,
      }
    % Slightly bigger margins than the latex defaults
    
    \geometry{verbose,tmargin=1in,bmargin=1in,lmargin=1in,rmargin=1in}
    
    

\begin{document}
    
    \maketitle
    
    

    
    \hypertarget{non-linear-modeling}{%
\section{Non-Linear Modeling}\label{non-linear-modeling}}

\href{https://mybinder.org/v2/gh/intro-stat-learning/ISLP_labs/v2.2?labpath=Ch07-nonlin-lab.ipynb}{\includesvg{https://mybinder.org/badge_logo.svg}}

    In this lab, we demonstrate some of the nonlinear models discussed in
this chapter. We use the \texttt{Wage} data as a running example, and
show that many of the complex non-linear fitting procedures discussed
can easily be implemented in \texttt{Python}.

As usual, we start with some of our standard imports.

    \begin{tcolorbox}[breakable, size=fbox, boxrule=1pt, pad at break*=1mm,colback=cellbackground, colframe=cellborder]
\prompt{In}{incolor}{1}{\boxspacing}
\begin{Verbatim}[commandchars=\\\{\}]
\PY{k+kn}{import}\PY{+w}{ }\PY{n+nn}{numpy}\PY{+w}{ }\PY{k}{as}\PY{+w}{ }\PY{n+nn}{np}\PY{o}{,}\PY{+w}{ }\PY{n+nn}{pandas}\PY{+w}{ }\PY{k}{as}\PY{+w}{ }\PY{n+nn}{pd}
\PY{k+kn}{from}\PY{+w}{ }\PY{n+nn}{matplotlib}\PY{n+nn}{.}\PY{n+nn}{pyplot}\PY{+w}{ }\PY{k+kn}{import} \PY{n}{subplots}
\PY{k+kn}{import}\PY{+w}{ }\PY{n+nn}{statsmodels}\PY{n+nn}{.}\PY{n+nn}{api}\PY{+w}{ }\PY{k}{as}\PY{+w}{ }\PY{n+nn}{sm}
\PY{k+kn}{from}\PY{+w}{ }\PY{n+nn}{ISLP}\PY{+w}{ }\PY{k+kn}{import} \PY{n}{load\PYZus{}data}
\PY{k+kn}{from}\PY{+w}{ }\PY{n+nn}{ISLP}\PY{n+nn}{.}\PY{n+nn}{models}\PY{+w}{ }\PY{k+kn}{import} \PY{p}{(}\PY{n}{summarize}\PY{p}{,}
                         \PY{n}{poly}\PY{p}{,}
                         \PY{n}{ModelSpec} \PY{k}{as} \PY{n}{MS}\PY{p}{)}
\PY{k+kn}{from}\PY{+w}{ }\PY{n+nn}{statsmodels}\PY{n+nn}{.}\PY{n+nn}{stats}\PY{n+nn}{.}\PY{n+nn}{anova}\PY{+w}{ }\PY{k+kn}{import} \PY{n}{anova\PYZus{}lm}
\end{Verbatim}
\end{tcolorbox}

    We again collect the new imports needed for this lab. Many of these are
developed specifically for the \texttt{ISLP} package.

    \begin{tcolorbox}[breakable, size=fbox, boxrule=1pt, pad at break*=1mm,colback=cellbackground, colframe=cellborder]
\prompt{In}{incolor}{2}{\boxspacing}
\begin{Verbatim}[commandchars=\\\{\}]
\PY{k+kn}{from}\PY{+w}{ }\PY{n+nn}{pygam}\PY{+w}{ }\PY{k+kn}{import} \PY{p}{(}\PY{n}{s} \PY{k}{as} \PY{n}{s\PYZus{}gam}\PY{p}{,}
                   \PY{n}{l} \PY{k}{as} \PY{n}{l\PYZus{}gam}\PY{p}{,}
                   \PY{n}{f} \PY{k}{as} \PY{n}{f\PYZus{}gam}\PY{p}{,}
                   \PY{n}{LinearGAM}\PY{p}{,}
                   \PY{n}{LogisticGAM}\PY{p}{)}

\PY{k+kn}{from}\PY{+w}{ }\PY{n+nn}{ISLP}\PY{n+nn}{.}\PY{n+nn}{transforms}\PY{+w}{ }\PY{k+kn}{import} \PY{p}{(}\PY{n}{BSpline}\PY{p}{,}
                             \PY{n}{NaturalSpline}\PY{p}{)}
\PY{k+kn}{from}\PY{+w}{ }\PY{n+nn}{ISLP}\PY{n+nn}{.}\PY{n+nn}{models}\PY{+w}{ }\PY{k+kn}{import} \PY{n}{bs}\PY{p}{,} \PY{n}{ns}
\PY{k+kn}{from}\PY{+w}{ }\PY{n+nn}{ISLP}\PY{n+nn}{.}\PY{n+nn}{pygam}\PY{+w}{ }\PY{k+kn}{import} \PY{p}{(}\PY{n}{approx\PYZus{}lam}\PY{p}{,}
                        \PY{n}{degrees\PYZus{}of\PYZus{}freedom}\PY{p}{,}
                        \PY{n}{plot} \PY{k}{as} \PY{n}{plot\PYZus{}gam}\PY{p}{,}
                        \PY{n}{anova} \PY{k}{as} \PY{n}{anova\PYZus{}gam}\PY{p}{)}
\end{Verbatim}
\end{tcolorbox}

    \hypertarget{polynomial-regression-and-step-functions}{%
\subsection{Polynomial Regression and Step
Functions}\label{polynomial-regression-and-step-functions}}

We start by demonstrating how Figure 7.1 can be reproduced. Let's begin
by loading the data.

    \begin{tcolorbox}[breakable, size=fbox, boxrule=1pt, pad at break*=1mm,colback=cellbackground, colframe=cellborder]
\prompt{In}{incolor}{3}{\boxspacing}
\begin{Verbatim}[commandchars=\\\{\}]
\PY{n}{Wage} \PY{o}{=} \PY{n}{load\PYZus{}data}\PY{p}{(}\PY{l+s+s1}{\PYZsq{}}\PY{l+s+s1}{Wage}\PY{l+s+s1}{\PYZsq{}}\PY{p}{)}
\PY{n}{y} \PY{o}{=} \PY{n}{Wage}\PY{p}{[}\PY{l+s+s1}{\PYZsq{}}\PY{l+s+s1}{wage}\PY{l+s+s1}{\PYZsq{}}\PY{p}{]}
\PY{n}{age} \PY{o}{=} \PY{n}{Wage}\PY{p}{[}\PY{l+s+s1}{\PYZsq{}}\PY{l+s+s1}{age}\PY{l+s+s1}{\PYZsq{}}\PY{p}{]}
\end{Verbatim}
\end{tcolorbox}

    Throughout most of this lab, our response is
\texttt{Wage{[}\textquotesingle{}wage\textquotesingle{}{]}}, which we
have stored as \texttt{y} above. As in Section 3.6.6, we will use the
\texttt{poly()} function to create a model matrix that will fit a
\(4\)th degree polynomial in \texttt{age}.

    \begin{tcolorbox}[breakable, size=fbox, boxrule=1pt, pad at break*=1mm,colback=cellbackground, colframe=cellborder]
\prompt{In}{incolor}{4}{\boxspacing}
\begin{Verbatim}[commandchars=\\\{\}]
\PY{n}{poly\PYZus{}age} \PY{o}{=} \PY{n}{MS}\PY{p}{(}\PY{p}{[}\PY{n}{poly}\PY{p}{(}\PY{l+s+s1}{\PYZsq{}}\PY{l+s+s1}{age}\PY{l+s+s1}{\PYZsq{}}\PY{p}{,} \PY{n}{degree}\PY{o}{=}\PY{l+m+mi}{4}\PY{p}{)}\PY{p}{]}\PY{p}{)}\PY{o}{.}\PY{n}{fit}\PY{p}{(}\PY{n}{Wage}\PY{p}{)}
\PY{n}{M} \PY{o}{=} \PY{n}{sm}\PY{o}{.}\PY{n}{OLS}\PY{p}{(}\PY{n}{y}\PY{p}{,} \PY{n}{poly\PYZus{}age}\PY{o}{.}\PY{n}{transform}\PY{p}{(}\PY{n}{Wage}\PY{p}{)}\PY{p}{)}\PY{o}{.}\PY{n}{fit}\PY{p}{(}\PY{p}{)}
\PY{n}{summarize}\PY{p}{(}\PY{n}{M}\PY{p}{)}
\end{Verbatim}
\end{tcolorbox}

            \begin{tcolorbox}[breakable, size=fbox, boxrule=.5pt, pad at break*=1mm, opacityfill=0]
\prompt{Out}{outcolor}{4}{\boxspacing}
\begin{Verbatim}[commandchars=\\\{\}]
                            coef  std err        t  P>|t|
intercept               111.7036    0.729  153.283  0.000
poly(age, degree=4)[0]  447.0679   39.915   11.201  0.000
poly(age, degree=4)[1] -478.3158   39.915  -11.983  0.000
poly(age, degree=4)[2]  125.5217   39.915    3.145  0.002
poly(age, degree=4)[3]  -77.9112   39.915   -1.952  0.051
\end{Verbatim}
\end{tcolorbox}
        
    This polynomial is constructed using the function \texttt{poly()}, which
creates a special \emph{transformer} \texttt{Poly()} (using
\texttt{sklearn} terminology for feature transformations such as
\texttt{PCA()} seen in Section 6.5.3) which allows for easy evaluation
of the polynomial at new data points. Here \texttt{poly()} is referred
to as a \emph{helper} function, and sets up the transformation;
\texttt{Poly()} is the actual workhorse that computes the
transformation. See also the discussion of transformations on page 118.

In the code above, the first line executes the \texttt{fit()} method
using the dataframe \texttt{Wage}. This recomputes and stores as
attributes any parameters needed by \texttt{Poly()} on the training
data, and these will be used on all subsequent evaluations of the
\texttt{transform()} method. For example, it is used on the second line,
as well as in the plotting function developed below.

    We now create a grid of values for \texttt{age} at which we want
predictions.

    \begin{tcolorbox}[breakable, size=fbox, boxrule=1pt, pad at break*=1mm,colback=cellbackground, colframe=cellborder]
\prompt{In}{incolor}{5}{\boxspacing}
\begin{Verbatim}[commandchars=\\\{\}]
\PY{n}{age\PYZus{}grid} \PY{o}{=} \PY{n}{np}\PY{o}{.}\PY{n}{linspace}\PY{p}{(}\PY{n}{age}\PY{o}{.}\PY{n}{min}\PY{p}{(}\PY{p}{)}\PY{p}{,}
                       \PY{n}{age}\PY{o}{.}\PY{n}{max}\PY{p}{(}\PY{p}{)}\PY{p}{,}
                       \PY{l+m+mi}{100}\PY{p}{)}
\PY{n}{age\PYZus{}df} \PY{o}{=} \PY{n}{pd}\PY{o}{.}\PY{n}{DataFrame}\PY{p}{(}\PY{p}{\PYZob{}}\PY{l+s+s1}{\PYZsq{}}\PY{l+s+s1}{age}\PY{l+s+s1}{\PYZsq{}}\PY{p}{:} \PY{n}{age\PYZus{}grid}\PY{p}{\PYZcb{}}\PY{p}{)}
\end{Verbatim}
\end{tcolorbox}

    Finally, we wish to plot the data and add the fit from the fourth-degree
polynomial. As we will make several similar plots below, we first write
a function to create all the ingredients and produce the plot. Our
function takes in a model specification (here a basis specified by a
transform), as well as a grid of \texttt{age} values. The function
produces a fitted curve as well as 95\% confidence bands. By using an
argument for \texttt{basis} we can produce and plot the results with
several different transforms, such as the splines we will see shortly.

    \begin{tcolorbox}[breakable, size=fbox, boxrule=1pt, pad at break*=1mm,colback=cellbackground, colframe=cellborder]
\prompt{In}{incolor}{6}{\boxspacing}
\begin{Verbatim}[commandchars=\\\{\}]
\PY{k}{def}\PY{+w}{ }\PY{n+nf}{plot\PYZus{}wage\PYZus{}fit}\PY{p}{(}\PY{n}{age\PYZus{}df}\PY{p}{,} 
                  \PY{n}{basis}\PY{p}{,}
                  \PY{n}{title}\PY{p}{)}\PY{p}{:}

    \PY{n}{X} \PY{o}{=} \PY{n}{basis}\PY{o}{.}\PY{n}{transform}\PY{p}{(}\PY{n}{Wage}\PY{p}{)}
    \PY{n}{Xnew} \PY{o}{=} \PY{n}{basis}\PY{o}{.}\PY{n}{transform}\PY{p}{(}\PY{n}{age\PYZus{}df}\PY{p}{)}
    \PY{n}{M} \PY{o}{=} \PY{n}{sm}\PY{o}{.}\PY{n}{OLS}\PY{p}{(}\PY{n}{y}\PY{p}{,} \PY{n}{X}\PY{p}{)}\PY{o}{.}\PY{n}{fit}\PY{p}{(}\PY{p}{)}
    \PY{n}{preds} \PY{o}{=} \PY{n}{M}\PY{o}{.}\PY{n}{get\PYZus{}prediction}\PY{p}{(}\PY{n}{Xnew}\PY{p}{)}
    \PY{n}{bands} \PY{o}{=} \PY{n}{preds}\PY{o}{.}\PY{n}{conf\PYZus{}int}\PY{p}{(}\PY{n}{alpha}\PY{o}{=}\PY{l+m+mf}{0.05}\PY{p}{)}
    \PY{n}{fig}\PY{p}{,} \PY{n}{ax} \PY{o}{=} \PY{n}{subplots}\PY{p}{(}\PY{n}{figsize}\PY{o}{=}\PY{p}{(}\PY{l+m+mi}{8}\PY{p}{,}\PY{l+m+mi}{8}\PY{p}{)}\PY{p}{)}
    \PY{n}{ax}\PY{o}{.}\PY{n}{scatter}\PY{p}{(}\PY{n}{age}\PY{p}{,}
               \PY{n}{y}\PY{p}{,}
               \PY{n}{facecolor}\PY{o}{=}\PY{l+s+s1}{\PYZsq{}}\PY{l+s+s1}{gray}\PY{l+s+s1}{\PYZsq{}}\PY{p}{,}
               \PY{n}{alpha}\PY{o}{=}\PY{l+m+mf}{0.5}\PY{p}{)}
    \PY{k}{for} \PY{n}{val}\PY{p}{,} \PY{n}{ls} \PY{o+ow}{in} \PY{n+nb}{zip}\PY{p}{(}\PY{p}{[}\PY{n}{preds}\PY{o}{.}\PY{n}{predicted\PYZus{}mean}\PY{p}{,}
                      \PY{n}{bands}\PY{p}{[}\PY{p}{:}\PY{p}{,}\PY{l+m+mi}{0}\PY{p}{]}\PY{p}{,}
                      \PY{n}{bands}\PY{p}{[}\PY{p}{:}\PY{p}{,}\PY{l+m+mi}{1}\PY{p}{]}\PY{p}{]}\PY{p}{,}
                     \PY{p}{[}\PY{l+s+s1}{\PYZsq{}}\PY{l+s+s1}{b}\PY{l+s+s1}{\PYZsq{}}\PY{p}{,}\PY{l+s+s1}{\PYZsq{}}\PY{l+s+s1}{r\PYZhy{}\PYZhy{}}\PY{l+s+s1}{\PYZsq{}}\PY{p}{,}\PY{l+s+s1}{\PYZsq{}}\PY{l+s+s1}{r\PYZhy{}\PYZhy{}}\PY{l+s+s1}{\PYZsq{}}\PY{p}{]}\PY{p}{)}\PY{p}{:}
        \PY{n}{ax}\PY{o}{.}\PY{n}{plot}\PY{p}{(}\PY{n}{age\PYZus{}df}\PY{o}{.}\PY{n}{values}\PY{p}{,} \PY{n}{val}\PY{p}{,} \PY{n}{ls}\PY{p}{,} \PY{n}{linewidth}\PY{o}{=}\PY{l+m+mi}{3}\PY{p}{)}
    \PY{n}{ax}\PY{o}{.}\PY{n}{set\PYZus{}title}\PY{p}{(}\PY{n}{title}\PY{p}{,} \PY{n}{fontsize}\PY{o}{=}\PY{l+m+mi}{20}\PY{p}{)}
    \PY{n}{ax}\PY{o}{.}\PY{n}{set\PYZus{}xlabel}\PY{p}{(}\PY{l+s+s1}{\PYZsq{}}\PY{l+s+s1}{Age}\PY{l+s+s1}{\PYZsq{}}\PY{p}{,} \PY{n}{fontsize}\PY{o}{=}\PY{l+m+mi}{20}\PY{p}{)}
    \PY{n}{ax}\PY{o}{.}\PY{n}{set\PYZus{}ylabel}\PY{p}{(}\PY{l+s+s1}{\PYZsq{}}\PY{l+s+s1}{Wage}\PY{l+s+s1}{\PYZsq{}}\PY{p}{,} \PY{n}{fontsize}\PY{o}{=}\PY{l+m+mi}{20}\PY{p}{)}\PY{p}{;}
    \PY{k}{return} \PY{n}{ax}
\end{Verbatim}
\end{tcolorbox}

    We include an argument \texttt{alpha} to \texttt{ax.scatter()} to add
some transparency to the points. This provides a visual indication of
density. Notice the use of the \texttt{zip()} function in the
\texttt{for} loop above (see Section 2.3.8). We have three lines to
plot, each with different colors and line types. Here \texttt{zip()}
conveniently bundles these together as iterators in the loop. \{In
\texttt{Python} speak, an ``iterator'' is an object with a finite number
of values, that can be iterated on, as in a loop.\}

We now plot the fit of the fourth-degree polynomial using this function.

    \begin{tcolorbox}[breakable, size=fbox, boxrule=1pt, pad at break*=1mm,colback=cellbackground, colframe=cellborder]
\prompt{In}{incolor}{7}{\boxspacing}
\begin{Verbatim}[commandchars=\\\{\}]
\PY{n}{plot\PYZus{}wage\PYZus{}fit}\PY{p}{(}\PY{n}{age\PYZus{}df}\PY{p}{,} 
              \PY{n}{poly\PYZus{}age}\PY{p}{,}
              \PY{l+s+s1}{\PYZsq{}}\PY{l+s+s1}{Degree\PYZhy{}4 Polynomial}\PY{l+s+s1}{\PYZsq{}}\PY{p}{)}\PY{p}{;}
\end{Verbatim}
\end{tcolorbox}

    \begin{center}
    \adjustimage{max size={0.9\linewidth}{0.9\paperheight}}{artifacts/latex/Ch07-nonlin-lab_files/artifacts/latex/Ch07-nonlin-lab_15_0.png}
    \end{center}
    { \hspace*{\fill} \\}
    
    

    With polynomial regression we must decide on the degree of the
polynomial to use. Sometimes we just wing it, and decide to use second
or third degree polynomials, simply to obtain a nonlinear fit. But we
can make such a decision in a more systematic way. One way to do this is
through hypothesis tests, which we demonstrate here. We now fit a series
of models ranging from linear (degree-one) to degree-five polynomials,
and look to determine the simplest model that is sufficient to explain
the relationship between \texttt{wage} and \texttt{age}. We use the
\texttt{anova\_lm()} function, which performs a series of ANOVA tests.
An \emph{analysis of
  variance} or ANOVA tests the null hypothesis that a model
\(\mathcal{M}_1\) is sufficient to explain the data against the
alternative hypothesis that a more complex model \(\mathcal{M}_2\) is
required. The determination is based on an F-test. To perform the test,
the models \(\mathcal{M}_1\) and \(\mathcal{M}_2\) must be
\emph{nested}: the space spanned by the predictors in \(\mathcal{M}_1\)
must be a subspace of the space spanned by the predictors in
\(\mathcal{M}_2\). In this case, we fit five different polynomial models
and sequentially compare the simpler model to the more complex model.

    \begin{tcolorbox}[breakable, size=fbox, boxrule=1pt, pad at break*=1mm,colback=cellbackground, colframe=cellborder]
\prompt{In}{incolor}{8}{\boxspacing}
\begin{Verbatim}[commandchars=\\\{\}]
\PY{n}{models} \PY{o}{=} \PY{p}{[}\PY{n}{MS}\PY{p}{(}\PY{p}{[}\PY{n}{poly}\PY{p}{(}\PY{l+s+s1}{\PYZsq{}}\PY{l+s+s1}{age}\PY{l+s+s1}{\PYZsq{}}\PY{p}{,} \PY{n}{degree}\PY{o}{=}\PY{n}{d}\PY{p}{)}\PY{p}{]}\PY{p}{)} 
          \PY{k}{for} \PY{n}{d} \PY{o+ow}{in} \PY{n+nb}{range}\PY{p}{(}\PY{l+m+mi}{1}\PY{p}{,} \PY{l+m+mi}{6}\PY{p}{)}\PY{p}{]}
\PY{n}{Xs} \PY{o}{=} \PY{p}{[}\PY{n}{model}\PY{o}{.}\PY{n}{fit\PYZus{}transform}\PY{p}{(}\PY{n}{Wage}\PY{p}{)} \PY{k}{for} \PY{n}{model} \PY{o+ow}{in} \PY{n}{models}\PY{p}{]}
\PY{n}{anova\PYZus{}lm}\PY{p}{(}\PY{o}{*}\PY{p}{[}\PY{n}{sm}\PY{o}{.}\PY{n}{OLS}\PY{p}{(}\PY{n}{y}\PY{p}{,} \PY{n}{X\PYZus{}}\PY{p}{)}\PY{o}{.}\PY{n}{fit}\PY{p}{(}\PY{p}{)}
           \PY{k}{for} \PY{n}{X\PYZus{}} \PY{o+ow}{in} \PY{n}{Xs}\PY{p}{]}\PY{p}{)}
\end{Verbatim}
\end{tcolorbox}

            \begin{tcolorbox}[breakable, size=fbox, boxrule=.5pt, pad at break*=1mm, opacityfill=0]
\prompt{Out}{outcolor}{8}{\boxspacing}
\begin{Verbatim}[commandchars=\\\{\}]
   df\_resid           ssr  df\_diff        ss\_diff           F        Pr(>F)
0    2998.0  5.022216e+06      0.0            NaN         NaN           NaN
1    2997.0  4.793430e+06      1.0  228786.010128  143.593107  2.363850e-32
2    2996.0  4.777674e+06      1.0   15755.693664    9.888756  1.679202e-03
3    2995.0  4.771604e+06      1.0    6070.152124    3.809813  5.104620e-02
4    2994.0  4.770322e+06      1.0    1282.563017    0.804976  3.696820e-01
\end{Verbatim}
\end{tcolorbox}
        
    Notice the \texttt{*} in the \texttt{anova\_lm()} line above. This
function takes a variable number of non-keyword arguments, in this case
fitted models. When these models are provided as a list (as is done
here), it must be prefixed by \texttt{*}.

The p-value comparing the linear \texttt{models{[}0{]}} to the quadratic
\texttt{models{[}1{]}} is essentially zero, indicating that a linear fit
is not sufficient. \{Indexing starting at zero is confusing for the
polynomial degree example, since \texttt{models{[}1{]}} is quadratic
rather than linear!\} Similarly the p-value comparing the quadratic
\texttt{models{[}1{]}} to the cubic \texttt{models{[}2{]}} is very low
(0.0017), so the quadratic fit is also insufficient. The p-value
comparing the cubic and degree-four polynomials, \texttt{models{[}2{]}}
and \texttt{models{[}3{]}}, is approximately 5\%, while the degree-five
polynomial \texttt{models{[}4{]}} seems unnecessary because its p-value
is 0.37. Hence, either a cubic or a quartic polynomial appear to provide
a reasonable fit to the data, but lower- or higher-order models are not
justified.

In this case, instead of using the \texttt{anova()} function, we could
have obtained these p-values more succinctly by exploiting the fact that
\texttt{poly()} creates orthogonal polynomials.

    \begin{tcolorbox}[breakable, size=fbox, boxrule=1pt, pad at break*=1mm,colback=cellbackground, colframe=cellborder]
\prompt{In}{incolor}{9}{\boxspacing}
\begin{Verbatim}[commandchars=\\\{\}]
\PY{n}{summarize}\PY{p}{(}\PY{n}{M}\PY{p}{)}
\end{Verbatim}
\end{tcolorbox}

            \begin{tcolorbox}[breakable, size=fbox, boxrule=.5pt, pad at break*=1mm, opacityfill=0]
\prompt{Out}{outcolor}{9}{\boxspacing}
\begin{Verbatim}[commandchars=\\\{\}]
                            coef  std err        t  P>|t|
intercept               111.7036    0.729  153.283  0.000
poly(age, degree=4)[0]  447.0679   39.915   11.201  0.000
poly(age, degree=4)[1] -478.3158   39.915  -11.983  0.000
poly(age, degree=4)[2]  125.5217   39.915    3.145  0.002
poly(age, degree=4)[3]  -77.9112   39.915   -1.952  0.051
\end{Verbatim}
\end{tcolorbox}
        
    Notice that the p-values are the same, and in fact the square of the
t-statistics are equal to the F-statistics from the \texttt{anova\_lm()}
function; for example:

    \begin{tcolorbox}[breakable, size=fbox, boxrule=1pt, pad at break*=1mm,colback=cellbackground, colframe=cellborder]
\prompt{In}{incolor}{10}{\boxspacing}
\begin{Verbatim}[commandchars=\\\{\}]
\PY{p}{(}\PY{o}{\PYZhy{}}\PY{l+m+mf}{11.983}\PY{p}{)}\PY{o}{*}\PY{o}{*}\PY{l+m+mi}{2}
\end{Verbatim}
\end{tcolorbox}

            \begin{tcolorbox}[breakable, size=fbox, boxrule=.5pt, pad at break*=1mm, opacityfill=0]
\prompt{Out}{outcolor}{10}{\boxspacing}
\begin{Verbatim}[commandchars=\\\{\}]
143.59228900000002
\end{Verbatim}
\end{tcolorbox}
        
    However, the ANOVA method works whether or not we used orthogonal
polynomials, provided the models are nested. For example, we can use
\texttt{anova\_lm()} to compare the following three models, which all
have a linear term in \texttt{education} and a polynomial in
\texttt{age} of different degrees:

    \begin{tcolorbox}[breakable, size=fbox, boxrule=1pt, pad at break*=1mm,colback=cellbackground, colframe=cellborder]
\prompt{In}{incolor}{11}{\boxspacing}
\begin{Verbatim}[commandchars=\\\{\}]
\PY{n}{models} \PY{o}{=} \PY{p}{[}\PY{n}{MS}\PY{p}{(}\PY{p}{[}\PY{l+s+s1}{\PYZsq{}}\PY{l+s+s1}{education}\PY{l+s+s1}{\PYZsq{}}\PY{p}{,} \PY{n}{poly}\PY{p}{(}\PY{l+s+s1}{\PYZsq{}}\PY{l+s+s1}{age}\PY{l+s+s1}{\PYZsq{}}\PY{p}{,} \PY{n}{degree}\PY{o}{=}\PY{n}{d}\PY{p}{)}\PY{p}{]}\PY{p}{)}
          \PY{k}{for} \PY{n}{d} \PY{o+ow}{in} \PY{n+nb}{range}\PY{p}{(}\PY{l+m+mi}{1}\PY{p}{,} \PY{l+m+mi}{4}\PY{p}{)}\PY{p}{]}
\PY{n}{XEs} \PY{o}{=} \PY{p}{[}\PY{n}{model}\PY{o}{.}\PY{n}{fit\PYZus{}transform}\PY{p}{(}\PY{n}{Wage}\PY{p}{)}
       \PY{k}{for} \PY{n}{model} \PY{o+ow}{in} \PY{n}{models}\PY{p}{]}
\PY{n}{anova\PYZus{}lm}\PY{p}{(}\PY{o}{*}\PY{p}{[}\PY{n}{sm}\PY{o}{.}\PY{n}{OLS}\PY{p}{(}\PY{n}{y}\PY{p}{,} \PY{n}{X\PYZus{}}\PY{p}{)}\PY{o}{.}\PY{n}{fit}\PY{p}{(}\PY{p}{)} \PY{k}{for} \PY{n}{X\PYZus{}} \PY{o+ow}{in} \PY{n}{XEs}\PY{p}{]}\PY{p}{)}
\end{Verbatim}
\end{tcolorbox}

            \begin{tcolorbox}[breakable, size=fbox, boxrule=.5pt, pad at break*=1mm, opacityfill=0]
\prompt{Out}{outcolor}{11}{\boxspacing}
\begin{Verbatim}[commandchars=\\\{\}]
   df\_resid           ssr  df\_diff        ss\_diff           F        Pr(>F)
0    2997.0  3.902335e+06      0.0            NaN         NaN           NaN
1    2996.0  3.759472e+06      1.0  142862.701185  113.991883  3.838075e-26
2    2995.0  3.753546e+06      1.0    5926.207070    4.728593  2.974318e-02
\end{Verbatim}
\end{tcolorbox}
        
    As an alternative to using hypothesis tests and ANOVA, we could choose
the polynomial degree using cross-validation, as discussed in Chapter 5.

Next we consider the task of predicting whether an individual earns more
than \$250,000 per year. We proceed much as before, except that first we
create the appropriate response vector, and then apply the
\texttt{glm()} function using the binomial family in order to fit a
polynomial logistic regression model.

    \begin{tcolorbox}[breakable, size=fbox, boxrule=1pt, pad at break*=1mm,colback=cellbackground, colframe=cellborder]
\prompt{In}{incolor}{12}{\boxspacing}
\begin{Verbatim}[commandchars=\\\{\}]
\PY{n}{X} \PY{o}{=} \PY{n}{poly\PYZus{}age}\PY{o}{.}\PY{n}{transform}\PY{p}{(}\PY{n}{Wage}\PY{p}{)}
\PY{n}{high\PYZus{}earn} \PY{o}{=} \PY{n}{Wage}\PY{p}{[}\PY{l+s+s1}{\PYZsq{}}\PY{l+s+s1}{high\PYZus{}earn}\PY{l+s+s1}{\PYZsq{}}\PY{p}{]} \PY{o}{=} \PY{n}{y} \PY{o}{\PYZgt{}} \PY{l+m+mi}{250} \PY{c+c1}{\PYZsh{} shorthand}
\PY{n}{glm} \PY{o}{=} \PY{n}{sm}\PY{o}{.}\PY{n}{GLM}\PY{p}{(}\PY{n}{y} \PY{o}{\PYZgt{}} \PY{l+m+mi}{250}\PY{p}{,}
             \PY{n}{X}\PY{p}{,}
             \PY{n}{family}\PY{o}{=}\PY{n}{sm}\PY{o}{.}\PY{n}{families}\PY{o}{.}\PY{n}{Binomial}\PY{p}{(}\PY{p}{)}\PY{p}{)}
\PY{n}{B} \PY{o}{=} \PY{n}{glm}\PY{o}{.}\PY{n}{fit}\PY{p}{(}\PY{p}{)}
\PY{n}{summarize}\PY{p}{(}\PY{n}{B}\PY{p}{)}
\end{Verbatim}
\end{tcolorbox}

            \begin{tcolorbox}[breakable, size=fbox, boxrule=.5pt, pad at break*=1mm, opacityfill=0]
\prompt{Out}{outcolor}{12}{\boxspacing}
\begin{Verbatim}[commandchars=\\\{\}]
                           coef  std err       z  P>|z|
intercept               -4.3012    0.345 -12.457  0.000
poly(age, degree=4)[0]  71.9642   26.133   2.754  0.006
poly(age, degree=4)[1] -85.7729   35.929  -2.387  0.017
poly(age, degree=4)[2]  34.1626   19.697   1.734  0.083
poly(age, degree=4)[3] -47.4008   24.105  -1.966  0.049
\end{Verbatim}
\end{tcolorbox}
        
    Once again, we make predictions using the \texttt{get\_prediction()}
method.

    \begin{tcolorbox}[breakable, size=fbox, boxrule=1pt, pad at break*=1mm,colback=cellbackground, colframe=cellborder]
\prompt{In}{incolor}{13}{\boxspacing}
\begin{Verbatim}[commandchars=\\\{\}]
\PY{n}{newX} \PY{o}{=} \PY{n}{poly\PYZus{}age}\PY{o}{.}\PY{n}{transform}\PY{p}{(}\PY{n}{age\PYZus{}df}\PY{p}{)}
\PY{n}{preds} \PY{o}{=} \PY{n}{B}\PY{o}{.}\PY{n}{get\PYZus{}prediction}\PY{p}{(}\PY{n}{newX}\PY{p}{)}
\PY{n}{bands} \PY{o}{=} \PY{n}{preds}\PY{o}{.}\PY{n}{conf\PYZus{}int}\PY{p}{(}\PY{n}{alpha}\PY{o}{=}\PY{l+m+mf}{0.05}\PY{p}{)}
\end{Verbatim}
\end{tcolorbox}

    We now plot the estimated relationship.

    \begin{tcolorbox}[breakable, size=fbox, boxrule=1pt, pad at break*=1mm,colback=cellbackground, colframe=cellborder]
\prompt{In}{incolor}{14}{\boxspacing}
\begin{Verbatim}[commandchars=\\\{\}]
\PY{n}{fig}\PY{p}{,} \PY{n}{ax} \PY{o}{=} \PY{n}{subplots}\PY{p}{(}\PY{n}{figsize}\PY{o}{=}\PY{p}{(}\PY{l+m+mi}{8}\PY{p}{,}\PY{l+m+mi}{8}\PY{p}{)}\PY{p}{)}
\PY{n}{rng} \PY{o}{=} \PY{n}{np}\PY{o}{.}\PY{n}{random}\PY{o}{.}\PY{n}{default\PYZus{}rng}\PY{p}{(}\PY{l+m+mi}{0}\PY{p}{)}
\PY{n}{ax}\PY{o}{.}\PY{n}{scatter}\PY{p}{(}\PY{n}{age} \PY{o}{+}
           \PY{l+m+mf}{0.2} \PY{o}{*} \PY{n}{rng}\PY{o}{.}\PY{n}{uniform}\PY{p}{(}\PY{n}{size}\PY{o}{=}\PY{n}{y}\PY{o}{.}\PY{n}{shape}\PY{p}{[}\PY{l+m+mi}{0}\PY{p}{]}\PY{p}{)}\PY{p}{,}
           \PY{n}{np}\PY{o}{.}\PY{n}{where}\PY{p}{(}\PY{n}{high\PYZus{}earn}\PY{p}{,} \PY{l+m+mf}{0.198}\PY{p}{,} \PY{l+m+mf}{0.002}\PY{p}{)}\PY{p}{,}
           \PY{n}{fc}\PY{o}{=}\PY{l+s+s1}{\PYZsq{}}\PY{l+s+s1}{gray}\PY{l+s+s1}{\PYZsq{}}\PY{p}{,}
           \PY{n}{marker}\PY{o}{=}\PY{l+s+s1}{\PYZsq{}}\PY{l+s+s1}{|}\PY{l+s+s1}{\PYZsq{}}\PY{p}{)}
\PY{k}{for} \PY{n}{val}\PY{p}{,} \PY{n}{ls} \PY{o+ow}{in} \PY{n+nb}{zip}\PY{p}{(}\PY{p}{[}\PY{n}{preds}\PY{o}{.}\PY{n}{predicted\PYZus{}mean}\PY{p}{,}
                  \PY{n}{bands}\PY{p}{[}\PY{p}{:}\PY{p}{,}\PY{l+m+mi}{0}\PY{p}{]}\PY{p}{,}
                  \PY{n}{bands}\PY{p}{[}\PY{p}{:}\PY{p}{,}\PY{l+m+mi}{1}\PY{p}{]}\PY{p}{]}\PY{p}{,}
                 \PY{p}{[}\PY{l+s+s1}{\PYZsq{}}\PY{l+s+s1}{b}\PY{l+s+s1}{\PYZsq{}}\PY{p}{,}\PY{l+s+s1}{\PYZsq{}}\PY{l+s+s1}{r\PYZhy{}\PYZhy{}}\PY{l+s+s1}{\PYZsq{}}\PY{p}{,}\PY{l+s+s1}{\PYZsq{}}\PY{l+s+s1}{r\PYZhy{}\PYZhy{}}\PY{l+s+s1}{\PYZsq{}}\PY{p}{]}\PY{p}{)}\PY{p}{:}
    \PY{n}{ax}\PY{o}{.}\PY{n}{plot}\PY{p}{(}\PY{n}{age\PYZus{}df}\PY{o}{.}\PY{n}{values}\PY{p}{,} \PY{n}{val}\PY{p}{,} \PY{n}{ls}\PY{p}{,} \PY{n}{linewidth}\PY{o}{=}\PY{l+m+mi}{3}\PY{p}{)}
\PY{n}{ax}\PY{o}{.}\PY{n}{set\PYZus{}title}\PY{p}{(}\PY{l+s+s1}{\PYZsq{}}\PY{l+s+s1}{Degree\PYZhy{}4 Polynomial}\PY{l+s+s1}{\PYZsq{}}\PY{p}{,} \PY{n}{fontsize}\PY{o}{=}\PY{l+m+mi}{20}\PY{p}{)}
\PY{n}{ax}\PY{o}{.}\PY{n}{set\PYZus{}xlabel}\PY{p}{(}\PY{l+s+s1}{\PYZsq{}}\PY{l+s+s1}{Age}\PY{l+s+s1}{\PYZsq{}}\PY{p}{,} \PY{n}{fontsize}\PY{o}{=}\PY{l+m+mi}{20}\PY{p}{)}
\PY{n}{ax}\PY{o}{.}\PY{n}{set\PYZus{}ylim}\PY{p}{(}\PY{p}{[}\PY{l+m+mi}{0}\PY{p}{,}\PY{l+m+mf}{0.2}\PY{p}{]}\PY{p}{)}
\PY{n}{ax}\PY{o}{.}\PY{n}{set\PYZus{}ylabel}\PY{p}{(}\PY{l+s+s1}{\PYZsq{}}\PY{l+s+s1}{P(Wage \PYZgt{} 250)}\PY{l+s+s1}{\PYZsq{}}\PY{p}{,} \PY{n}{fontsize}\PY{o}{=}\PY{l+m+mi}{20}\PY{p}{)}\PY{p}{;}
\end{Verbatim}
\end{tcolorbox}

    \begin{center}
    \adjustimage{max size={0.9\linewidth}{0.9\paperheight}}{artifacts/latex/Ch07-nonlin-lab_files/artifacts/latex/Ch07-nonlin-lab_30_0.png}
    \end{center}
    { \hspace*{\fill} \\}
    
    We have drawn the \texttt{age} values corresponding to the observations
with \texttt{wage} values above 250 as gray marks on the top of the
plot, and those with \texttt{wage} values below 250 are shown as gray
marks on the bottom of the plot. We added a small amount of noise to
jitter the \texttt{age} values a bit so that observations with the same
\texttt{age} value do not cover each other up. This type of plot is
often called a \emph{rug plot}.

In order to fit a step function, as discussed in Section 7.2, we first
use the \texttt{pd.qcut()} function to discretize \texttt{age} based on
quantiles. Then we use \texttt{pd.get\_dummies()} to create the columns
of the model matrix for this categorical variable. Note that this
function will include \emph{all} columns for a given categorical, rather
than the usual approach which drops one of the levels.

    \begin{tcolorbox}[breakable, size=fbox, boxrule=1pt, pad at break*=1mm,colback=cellbackground, colframe=cellborder]
\prompt{In}{incolor}{15}{\boxspacing}
\begin{Verbatim}[commandchars=\\\{\}]
\PY{n}{cut\PYZus{}age} \PY{o}{=} \PY{n}{pd}\PY{o}{.}\PY{n}{qcut}\PY{p}{(}\PY{n}{age}\PY{p}{,} \PY{l+m+mi}{4}\PY{p}{)}
\PY{n}{summarize}\PY{p}{(}\PY{n}{sm}\PY{o}{.}\PY{n}{OLS}\PY{p}{(}\PY{n}{y}\PY{p}{,} \PY{n}{pd}\PY{o}{.}\PY{n}{get\PYZus{}dummies}\PY{p}{(}\PY{n}{cut\PYZus{}age}\PY{p}{)}\PY{p}{)}\PY{o}{.}\PY{n}{fit}\PY{p}{(}\PY{p}{)}\PY{p}{)}
\end{Verbatim}
\end{tcolorbox}

            \begin{tcolorbox}[breakable, size=fbox, boxrule=.5pt, pad at break*=1mm, opacityfill=0]
\prompt{Out}{outcolor}{15}{\boxspacing}
\begin{Verbatim}[commandchars=\\\{\}]
                     coef  std err       t  P>|t|
(17.999, 33.75]   94.1584    1.478  63.692    0.0
(33.75, 42.0]    116.6608    1.470  79.385    0.0
(42.0, 51.0]     119.1887    1.416  84.147    0.0
(51.0, 80.0]     116.5717    1.559  74.751    0.0
\end{Verbatim}
\end{tcolorbox}
        
    Here \texttt{pd.qcut()} automatically picked the cutpoints based on the
quantiles 25\%, 50\% and 75\%, which results in four regions. We could
also have specified our own quantiles directly instead of the argument
\texttt{4}. For cuts not based on quantiles we would use the
\texttt{pd.cut()} function. The function \texttt{pd.qcut()} (and
\texttt{pd.cut()}) returns an ordered categorical variable. The
regression model then creates a set of dummy variables for use in the
regression. Since \texttt{age} is the only variable in the model, the
value \$94,158.40 is the average salary for those under 33.75 years of
age, and the other coefficients are the average salary for those in the
other age groups. We can produce predictions and plots just as we did in
the case of the polynomial fit.

    \hypertarget{splines}{%
\subsection{Splines}\label{splines}}

In order to fit regression splines, we use transforms from the
\texttt{ISLP} package. The actual spline evaluation functions are in the
\texttt{scipy.interpolate} package; we have simply wrapped them as
transforms similar to \texttt{Poly()} and \texttt{PCA()}.

In Section 7.4, we saw that regression splines can be fit by
constructing an appropriate matrix of basis functions. The
\texttt{BSpline()} function generates the entire matrix of basis
functions for splines with the specified set of knots. By default, the
B-splines produced are cubic. To change the degree, use the argument
\texttt{degree}.

    \begin{tcolorbox}[breakable, size=fbox, boxrule=1pt, pad at break*=1mm,colback=cellbackground, colframe=cellborder]
\prompt{In}{incolor}{16}{\boxspacing}
\begin{Verbatim}[commandchars=\\\{\}]
\PY{n}{bs\PYZus{}} \PY{o}{=} \PY{n}{BSpline}\PY{p}{(}\PY{n}{internal\PYZus{}knots}\PY{o}{=}\PY{p}{[}\PY{l+m+mi}{25}\PY{p}{,}\PY{l+m+mi}{40}\PY{p}{,}\PY{l+m+mi}{60}\PY{p}{]}\PY{p}{,} \PY{n}{intercept}\PY{o}{=}\PY{k+kc}{True}\PY{p}{)}\PY{o}{.}\PY{n}{fit}\PY{p}{(}\PY{n}{age}\PY{p}{)}
\PY{n}{bs\PYZus{}age} \PY{o}{=} \PY{n}{bs\PYZus{}}\PY{o}{.}\PY{n}{transform}\PY{p}{(}\PY{n}{age}\PY{p}{)}
\PY{n}{bs\PYZus{}age}\PY{o}{.}\PY{n}{shape}
\end{Verbatim}
\end{tcolorbox}

            \begin{tcolorbox}[breakable, size=fbox, boxrule=.5pt, pad at break*=1mm, opacityfill=0]
\prompt{Out}{outcolor}{16}{\boxspacing}
\begin{Verbatim}[commandchars=\\\{\}]
(3000, 7)
\end{Verbatim}
\end{tcolorbox}
        
    This results in a seven-column matrix, which is what is expected for a
cubic-spline basis with 3 interior knots. We can form this same matrix
using the \texttt{bs()} object, which facilitates adding this to a
model-matrix builder (as in \texttt{poly()} versus its workhorse
\texttt{Poly()}) described in Section 7.8.1.

We now fit a cubic spline model to the \texttt{Wage} data.

    \begin{tcolorbox}[breakable, size=fbox, boxrule=1pt, pad at break*=1mm,colback=cellbackground, colframe=cellborder]
\prompt{In}{incolor}{17}{\boxspacing}
\begin{Verbatim}[commandchars=\\\{\}]
\PY{n}{bs\PYZus{}age} \PY{o}{=} \PY{n}{MS}\PY{p}{(}\PY{p}{[}\PY{n}{bs}\PY{p}{(}\PY{l+s+s1}{\PYZsq{}}\PY{l+s+s1}{age}\PY{l+s+s1}{\PYZsq{}}\PY{p}{,} \PY{n}{internal\PYZus{}knots}\PY{o}{=}\PY{p}{[}\PY{l+m+mi}{25}\PY{p}{,}\PY{l+m+mi}{40}\PY{p}{,}\PY{l+m+mi}{60}\PY{p}{]}\PY{p}{)}\PY{p}{]}\PY{p}{)}
\PY{n}{Xbs} \PY{o}{=} \PY{n}{bs\PYZus{}age}\PY{o}{.}\PY{n}{fit\PYZus{}transform}\PY{p}{(}\PY{n}{Wage}\PY{p}{)}
\PY{n}{M} \PY{o}{=} \PY{n}{sm}\PY{o}{.}\PY{n}{OLS}\PY{p}{(}\PY{n}{y}\PY{p}{,} \PY{n}{Xbs}\PY{p}{)}\PY{o}{.}\PY{n}{fit}\PY{p}{(}\PY{p}{)}
\PY{n}{summarize}\PY{p}{(}\PY{n}{M}\PY{p}{)}
\end{Verbatim}
\end{tcolorbox}

            \begin{tcolorbox}[breakable, size=fbox, boxrule=.5pt, pad at break*=1mm, opacityfill=0]
\prompt{Out}{outcolor}{17}{\boxspacing}
\begin{Verbatim}[commandchars=\\\{\}]
                                            coef  std err      t  P>|t|
intercept                                60.4937    9.460  6.394  0.000
bs(age, internal\_knots=[25, 40, 60])[0]   3.9805   12.538  0.317  0.751
bs(age, internal\_knots=[25, 40, 60])[1]  44.6310    9.626  4.636  0.000
bs(age, internal\_knots=[25, 40, 60])[2]  62.8388   10.755  5.843  0.000
bs(age, internal\_knots=[25, 40, 60])[3]  55.9908   10.706  5.230  0.000
bs(age, internal\_knots=[25, 40, 60])[4]  50.6881   14.402  3.520  0.000
bs(age, internal\_knots=[25, 40, 60])[5]  16.6061   19.126  0.868  0.385
\end{Verbatim}
\end{tcolorbox}
        
    The column names are a little cumbersome, and have caused us to truncate
the printed summary. They can be set on construction using the
\texttt{name} argument as follows.

    \begin{tcolorbox}[breakable, size=fbox, boxrule=1pt, pad at break*=1mm,colback=cellbackground, colframe=cellborder]
\prompt{In}{incolor}{18}{\boxspacing}
\begin{Verbatim}[commandchars=\\\{\}]
\PY{n}{bs\PYZus{}age} \PY{o}{=} \PY{n}{MS}\PY{p}{(}\PY{p}{[}\PY{n}{bs}\PY{p}{(}\PY{l+s+s1}{\PYZsq{}}\PY{l+s+s1}{age}\PY{l+s+s1}{\PYZsq{}}\PY{p}{,}
                \PY{n}{internal\PYZus{}knots}\PY{o}{=}\PY{p}{[}\PY{l+m+mi}{25}\PY{p}{,}\PY{l+m+mi}{40}\PY{p}{,}\PY{l+m+mi}{60}\PY{p}{]}\PY{p}{,}
                \PY{n}{name}\PY{o}{=}\PY{l+s+s1}{\PYZsq{}}\PY{l+s+s1}{bs(age)}\PY{l+s+s1}{\PYZsq{}}\PY{p}{)}\PY{p}{]}\PY{p}{)}
\PY{n}{Xbs} \PY{o}{=} \PY{n}{bs\PYZus{}age}\PY{o}{.}\PY{n}{fit\PYZus{}transform}\PY{p}{(}\PY{n}{Wage}\PY{p}{)}
\PY{n}{M} \PY{o}{=} \PY{n}{sm}\PY{o}{.}\PY{n}{OLS}\PY{p}{(}\PY{n}{y}\PY{p}{,} \PY{n}{Xbs}\PY{p}{)}\PY{o}{.}\PY{n}{fit}\PY{p}{(}\PY{p}{)}
\PY{n}{summarize}\PY{p}{(}\PY{n}{M}\PY{p}{)}
\end{Verbatim}
\end{tcolorbox}

            \begin{tcolorbox}[breakable, size=fbox, boxrule=.5pt, pad at break*=1mm, opacityfill=0]
\prompt{Out}{outcolor}{18}{\boxspacing}
\begin{Verbatim}[commandchars=\\\{\}]
               coef  std err      t  P>|t|
intercept   60.4937    9.460  6.394  0.000
bs(age)[0]   3.9805   12.538  0.317  0.751
bs(age)[1]  44.6310    9.626  4.636  0.000
bs(age)[2]  62.8388   10.755  5.843  0.000
bs(age)[3]  55.9908   10.706  5.230  0.000
bs(age)[4]  50.6881   14.402  3.520  0.000
bs(age)[5]  16.6061   19.126  0.868  0.385
\end{Verbatim}
\end{tcolorbox}
        
    Notice that there are 6 spline coefficients rather than 7. This is
because, by default, \texttt{bs()} assumes \texttt{intercept=False},
since we typically have an overall intercept in the model. So it
generates the spline basis with the given knots, and then discards one
of the basis functions to account for the intercept.

We could also use the \texttt{df} (degrees of freedom) option to specify
the complexity of the spline. We see above that with 3 knots, the spline
basis has 6 columns or degrees of freedom. When we specify \texttt{df=6}
rather than the actual knots, \texttt{bs()} will produce a spline with 3
knots chosen at uniform quantiles of the training data. We can see these
chosen knots most easily using \texttt{Bspline()} directly:

    \begin{tcolorbox}[breakable, size=fbox, boxrule=1pt, pad at break*=1mm,colback=cellbackground, colframe=cellborder]
\prompt{In}{incolor}{19}{\boxspacing}
\begin{Verbatim}[commandchars=\\\{\}]
\PY{n}{BSpline}\PY{p}{(}\PY{n}{df}\PY{o}{=}\PY{l+m+mi}{6}\PY{p}{)}\PY{o}{.}\PY{n}{fit}\PY{p}{(}\PY{n}{age}\PY{p}{)}\PY{o}{.}\PY{n}{internal\PYZus{}knots\PYZus{}}
\end{Verbatim}
\end{tcolorbox}

            \begin{tcolorbox}[breakable, size=fbox, boxrule=.5pt, pad at break*=1mm, opacityfill=0]
\prompt{Out}{outcolor}{19}{\boxspacing}
\begin{Verbatim}[commandchars=\\\{\}]
array([33.75, 42.  , 51.  ])
\end{Verbatim}
\end{tcolorbox}
        
    When asking for six degrees of freedom, the transform chooses knots at
ages 33.75, 42.0, and 51.0, which correspond to the 25th, 50th, and 75th
percentiles of \texttt{age}.

When using B-splines we need not limit ourselves to cubic polynomials
(i.e.~\texttt{degree=3}). For instance, using \texttt{degree=0} results
in piecewise constant functions, as in our example with
\texttt{pd.qcut()} above.

    \begin{tcolorbox}[breakable, size=fbox, boxrule=1pt, pad at break*=1mm,colback=cellbackground, colframe=cellborder]
\prompt{In}{incolor}{20}{\boxspacing}
\begin{Verbatim}[commandchars=\\\{\}]
\PY{n}{bs\PYZus{}age0} \PY{o}{=} \PY{n}{MS}\PY{p}{(}\PY{p}{[}\PY{n}{bs}\PY{p}{(}\PY{l+s+s1}{\PYZsq{}}\PY{l+s+s1}{age}\PY{l+s+s1}{\PYZsq{}}\PY{p}{,}
                 \PY{n}{df}\PY{o}{=}\PY{l+m+mi}{3}\PY{p}{,} 
                 \PY{n}{degree}\PY{o}{=}\PY{l+m+mi}{0}\PY{p}{)}\PY{p}{]}\PY{p}{)}\PY{o}{.}\PY{n}{fit}\PY{p}{(}\PY{n}{Wage}\PY{p}{)}
\PY{n}{Xbs0} \PY{o}{=} \PY{n}{bs\PYZus{}age0}\PY{o}{.}\PY{n}{transform}\PY{p}{(}\PY{n}{Wage}\PY{p}{)}
\PY{n}{summarize}\PY{p}{(}\PY{n}{sm}\PY{o}{.}\PY{n}{OLS}\PY{p}{(}\PY{n}{y}\PY{p}{,} \PY{n}{Xbs0}\PY{p}{)}\PY{o}{.}\PY{n}{fit}\PY{p}{(}\PY{p}{)}\PY{p}{)}
\end{Verbatim}
\end{tcolorbox}

            \begin{tcolorbox}[breakable, size=fbox, boxrule=.5pt, pad at break*=1mm, opacityfill=0]
\prompt{Out}{outcolor}{20}{\boxspacing}
\begin{Verbatim}[commandchars=\\\{\}]
                               coef  std err       t  P>|t|
intercept                   94.1584    1.478  63.687    0.0
bs(age, df=3, degree=0)[0]  22.3490    2.152  10.388    0.0
bs(age, df=3, degree=0)[1]  24.8076    2.044  12.137    0.0
bs(age, df=3, degree=0)[2]  22.7814    2.087  10.917    0.0
\end{Verbatim}
\end{tcolorbox}
        
    This fit should be compared with cell {[}15{]} where we use
\texttt{qcut()} to create four bins by cutting at the 25\%, 50\% and
75\% quantiles of \texttt{age}. Since we specified \texttt{df=3} for
degree-zero splines here, there will also be knots at the same three
quantiles. Although the coefficients appear different, we see that this
is a result of the different coding. For example, the first coefficient
is identical in both cases, and is the mean response in the first bin.
For the second coefficient, we have
\(94.158 + 22.349 = 116.507 \approx 116.611\), the latter being the mean
in the second bin in cell {[}15{]}. Here the intercept is coded by a
column of ones, so the second, third and fourth coefficients are
increments for those bins. Why is the sum not exactly the same? It turns
out that the \texttt{qcut()} uses \(\leq\), while \texttt{bs()} uses
\(<\) when deciding bin membership.

    In order to fit a natural spline, we use the \texttt{NaturalSpline()}
transform with the corresponding helper \texttt{ns()}. Here we fit a
natural spline with five degrees of freedom (excluding the intercept)
and plot the results.

    \begin{tcolorbox}[breakable, size=fbox, boxrule=1pt, pad at break*=1mm,colback=cellbackground, colframe=cellborder]
\prompt{In}{incolor}{21}{\boxspacing}
\begin{Verbatim}[commandchars=\\\{\}]
\PY{n}{ns\PYZus{}age} \PY{o}{=} \PY{n}{MS}\PY{p}{(}\PY{p}{[}\PY{n}{ns}\PY{p}{(}\PY{l+s+s1}{\PYZsq{}}\PY{l+s+s1}{age}\PY{l+s+s1}{\PYZsq{}}\PY{p}{,} \PY{n}{df}\PY{o}{=}\PY{l+m+mi}{5}\PY{p}{)}\PY{p}{]}\PY{p}{)}\PY{o}{.}\PY{n}{fit}\PY{p}{(}\PY{n}{Wage}\PY{p}{)}
\PY{n}{M\PYZus{}ns} \PY{o}{=} \PY{n}{sm}\PY{o}{.}\PY{n}{OLS}\PY{p}{(}\PY{n}{y}\PY{p}{,} \PY{n}{ns\PYZus{}age}\PY{o}{.}\PY{n}{transform}\PY{p}{(}\PY{n}{Wage}\PY{p}{)}\PY{p}{)}\PY{o}{.}\PY{n}{fit}\PY{p}{(}\PY{p}{)}
\PY{n}{summarize}\PY{p}{(}\PY{n}{M\PYZus{}ns}\PY{p}{)}
\end{Verbatim}
\end{tcolorbox}

            \begin{tcolorbox}[breakable, size=fbox, boxrule=.5pt, pad at break*=1mm, opacityfill=0]
\prompt{Out}{outcolor}{21}{\boxspacing}
\begin{Verbatim}[commandchars=\\\{\}]
                     coef  std err       t  P>|t|
intercept         60.4752    4.708  12.844  0.000
ns(age, df=5)[0]  61.5267    4.709  13.065  0.000
ns(age, df=5)[1]  55.6912    5.717   9.741  0.000
ns(age, df=5)[2]  46.8184    4.948   9.463  0.000
ns(age, df=5)[3]  83.2036   11.918   6.982  0.000
ns(age, df=5)[4]   6.8770    9.484   0.725  0.468
\end{Verbatim}
\end{tcolorbox}
        
    We now plot the natural spline using our plotting function.

    \begin{tcolorbox}[breakable, size=fbox, boxrule=1pt, pad at break*=1mm,colback=cellbackground, colframe=cellborder]
\prompt{In}{incolor}{22}{\boxspacing}
\begin{Verbatim}[commandchars=\\\{\}]
\PY{n}{plot\PYZus{}wage\PYZus{}fit}\PY{p}{(}\PY{n}{age\PYZus{}df}\PY{p}{,}
              \PY{n}{ns\PYZus{}age}\PY{p}{,}
              \PY{l+s+s1}{\PYZsq{}}\PY{l+s+s1}{Natural spline, df=5}\PY{l+s+s1}{\PYZsq{}}\PY{p}{)}\PY{p}{;}
\end{Verbatim}
\end{tcolorbox}

    \begin{center}
    \adjustimage{max size={0.9\linewidth}{0.9\paperheight}}{artifacts/latex/Ch07-nonlin-lab_files/artifacts/latex/Ch07-nonlin-lab_48_0.png}
    \end{center}
    { \hspace*{\fill} \\}
    
    \hypertarget{smoothing-splines-and-gams}{%
\subsection{Smoothing Splines and
GAMs}\label{smoothing-splines-and-gams}}

A smoothing spline is a special case of a GAM with squared-error loss
and a single feature. To fit GAMs in \texttt{Python} we will use the
\texttt{pygam} package which can be installed via
\texttt{pip\ install\ pygam}. The estimator \texttt{LinearGAM()} uses
squared-error loss. The GAM is specified by associating each column of a
model matrix with a particular smoothing operation: \texttt{s} for
smoothing spline; \texttt{l} for linear, and \texttt{f} for factor or
categorical variables. The argument \texttt{0} passed to \texttt{s}
below indicates that this smoother will apply to the first column of a
feature matrix. Below, we pass it a matrix with a single column:
\texttt{X\_age}. The argument \texttt{lam} is the penalty parameter
\(\lambda\) as discussed in Section 7.5.2.

    \begin{tcolorbox}[breakable, size=fbox, boxrule=1pt, pad at break*=1mm,colback=cellbackground, colframe=cellborder]
\prompt{In}{incolor}{23}{\boxspacing}
\begin{Verbatim}[commandchars=\\\{\}]
\PY{n}{X\PYZus{}age} \PY{o}{=} \PY{n}{np}\PY{o}{.}\PY{n}{asarray}\PY{p}{(}\PY{n}{age}\PY{p}{)}\PY{o}{.}\PY{n}{reshape}\PY{p}{(}\PY{p}{(}\PY{o}{\PYZhy{}}\PY{l+m+mi}{1}\PY{p}{,}\PY{l+m+mi}{1}\PY{p}{)}\PY{p}{)}
\PY{n}{gam} \PY{o}{=} \PY{n}{LinearGAM}\PY{p}{(}\PY{n}{s\PYZus{}gam}\PY{p}{(}\PY{l+m+mi}{0}\PY{p}{,} \PY{n}{lam}\PY{o}{=}\PY{l+m+mf}{0.6}\PY{p}{)}\PY{p}{)}
\PY{n}{gam}\PY{o}{.}\PY{n}{fit}\PY{p}{(}\PY{n}{X\PYZus{}age}\PY{p}{,} \PY{n}{y}\PY{p}{)}
\end{Verbatim}
\end{tcolorbox}

            \begin{tcolorbox}[breakable, size=fbox, boxrule=.5pt, pad at break*=1mm, opacityfill=0]
\prompt{Out}{outcolor}{23}{\boxspacing}
\begin{Verbatim}[commandchars=\\\{\}]
LinearGAM(callbacks=[Deviance(), Diffs()], fit\_intercept=True,
   max\_iter=100, scale=None, terms=s(0) + intercept, tol=0.0001,
   verbose=False)
\end{Verbatim}
\end{tcolorbox}
        
    The \texttt{pygam} library generally expects a matrix of features so we
reshape \texttt{age} to be a matrix (a two-dimensional array) instead of
a vector (i.e.~a one-dimensional array). The \texttt{-1} in the call to
the \texttt{reshape()} method tells \texttt{numpy} to impute the size of
that dimension based on the remaining entries of the shape tuple.

Let's investigate how the fit changes with the smoothing parameter
\texttt{lam}. The function \texttt{np.logspace()} is similar to
\texttt{np.linspace()} but spaces points evenly on the log-scale. Below
we vary \texttt{lam} from \(10^{-2}\) to \(10^6\).

    \begin{tcolorbox}[breakable, size=fbox, boxrule=1pt, pad at break*=1mm,colback=cellbackground, colframe=cellborder]
\prompt{In}{incolor}{24}{\boxspacing}
\begin{Verbatim}[commandchars=\\\{\}]
\PY{n}{fig}\PY{p}{,} \PY{n}{ax} \PY{o}{=} \PY{n}{subplots}\PY{p}{(}\PY{n}{figsize}\PY{o}{=}\PY{p}{(}\PY{l+m+mi}{8}\PY{p}{,}\PY{l+m+mi}{8}\PY{p}{)}\PY{p}{)}
\PY{n}{ax}\PY{o}{.}\PY{n}{scatter}\PY{p}{(}\PY{n}{age}\PY{p}{,} \PY{n}{y}\PY{p}{,} \PY{n}{facecolor}\PY{o}{=}\PY{l+s+s1}{\PYZsq{}}\PY{l+s+s1}{gray}\PY{l+s+s1}{\PYZsq{}}\PY{p}{,} \PY{n}{alpha}\PY{o}{=}\PY{l+m+mf}{0.5}\PY{p}{)}
\PY{k}{for} \PY{n}{lam} \PY{o+ow}{in} \PY{n}{np}\PY{o}{.}\PY{n}{logspace}\PY{p}{(}\PY{o}{\PYZhy{}}\PY{l+m+mi}{2}\PY{p}{,} \PY{l+m+mi}{6}\PY{p}{,} \PY{l+m+mi}{5}\PY{p}{)}\PY{p}{:}
    \PY{n}{gam} \PY{o}{=} \PY{n}{LinearGAM}\PY{p}{(}\PY{n}{s\PYZus{}gam}\PY{p}{(}\PY{l+m+mi}{0}\PY{p}{,} \PY{n}{lam}\PY{o}{=}\PY{n}{lam}\PY{p}{)}\PY{p}{)}\PY{o}{.}\PY{n}{fit}\PY{p}{(}\PY{n}{X\PYZus{}age}\PY{p}{,} \PY{n}{y}\PY{p}{)}
    \PY{n}{ax}\PY{o}{.}\PY{n}{plot}\PY{p}{(}\PY{n}{age\PYZus{}grid}\PY{p}{,}
            \PY{n}{gam}\PY{o}{.}\PY{n}{predict}\PY{p}{(}\PY{n}{age\PYZus{}grid}\PY{p}{)}\PY{p}{,}
            \PY{n}{label}\PY{o}{=}\PY{l+s+s1}{\PYZsq{}}\PY{l+s+si}{\PYZob{}:.1e\PYZcb{}}\PY{l+s+s1}{\PYZsq{}}\PY{o}{.}\PY{n}{format}\PY{p}{(}\PY{n}{lam}\PY{p}{)}\PY{p}{,}
            \PY{n}{linewidth}\PY{o}{=}\PY{l+m+mi}{3}\PY{p}{)}
\PY{n}{ax}\PY{o}{.}\PY{n}{set\PYZus{}xlabel}\PY{p}{(}\PY{l+s+s1}{\PYZsq{}}\PY{l+s+s1}{Age}\PY{l+s+s1}{\PYZsq{}}\PY{p}{,} \PY{n}{fontsize}\PY{o}{=}\PY{l+m+mi}{20}\PY{p}{)}
\PY{n}{ax}\PY{o}{.}\PY{n}{set\PYZus{}ylabel}\PY{p}{(}\PY{l+s+s1}{\PYZsq{}}\PY{l+s+s1}{Wage}\PY{l+s+s1}{\PYZsq{}}\PY{p}{,} \PY{n}{fontsize}\PY{o}{=}\PY{l+m+mi}{20}\PY{p}{)}\PY{p}{;}
\PY{n}{ax}\PY{o}{.}\PY{n}{legend}\PY{p}{(}\PY{n}{title}\PY{o}{=}\PY{l+s+s1}{\PYZsq{}}\PY{l+s+s1}{\PYZdl{}}\PY{l+s+s1}{\PYZbs{}}\PY{l+s+s1}{lambda\PYZdl{}}\PY{l+s+s1}{\PYZsq{}}\PY{p}{)}\PY{p}{;}
\end{Verbatim}
\end{tcolorbox}

    \begin{center}
    \adjustimage{max size={0.9\linewidth}{0.9\paperheight}}{artifacts/latex/Ch07-nonlin-lab_files/artifacts/latex/Ch07-nonlin-lab_52_0.png}
    \end{center}
    { \hspace*{\fill} \\}
    
    The \texttt{pygam} package can perform a search for an optimal smoothing
parameter.

    \begin{tcolorbox}[breakable, size=fbox, boxrule=1pt, pad at break*=1mm,colback=cellbackground, colframe=cellborder]
\prompt{In}{incolor}{25}{\boxspacing}
\begin{Verbatim}[commandchars=\\\{\}]
\PY{n}{gam\PYZus{}opt} \PY{o}{=} \PY{n}{gam}\PY{o}{.}\PY{n}{gridsearch}\PY{p}{(}\PY{n}{X\PYZus{}age}\PY{p}{,} \PY{n}{y}\PY{p}{)}
\PY{n}{ax}\PY{o}{.}\PY{n}{plot}\PY{p}{(}\PY{n}{age\PYZus{}grid}\PY{p}{,}
        \PY{n}{gam\PYZus{}opt}\PY{o}{.}\PY{n}{predict}\PY{p}{(}\PY{n}{age\PYZus{}grid}\PY{p}{)}\PY{p}{,}
        \PY{n}{label}\PY{o}{=}\PY{l+s+s1}{\PYZsq{}}\PY{l+s+s1}{Grid search}\PY{l+s+s1}{\PYZsq{}}\PY{p}{,}
        \PY{n}{linewidth}\PY{o}{=}\PY{l+m+mi}{4}\PY{p}{)}
\PY{n}{ax}\PY{o}{.}\PY{n}{legend}\PY{p}{(}\PY{p}{)}
\PY{n}{fig}
\end{Verbatim}
\end{tcolorbox}

    \begin{Verbatim}[commandchars=\\\{\}]
 \def\tcRGB{\textcolor[RGB]}\expandafter\tcRGB\expandafter{\detokenize{255,0,0}}{  0\%} \def\tcRGB{\textcolor[RGB]}\expandafter\tcRGB\expandafter{\detokenize{255,0,0}}{(0 of 11)} |
| Elapsed Time: 0:00:00 ETA:  --:--:--
    \end{Verbatim}

    \begin{Verbatim}[commandchars=\\\{\}]
 \def\tcRGB{\textcolor[RGB]}\expandafter\tcRGB\expandafter{\detokenize{255,105,0}}{ 18\%} \def\tcRGB{\textcolor[RGB]}\expandafter\tcRGB\expandafter{\detokenize{255,105,0}}{(2 of 11)} |\#\#\#\#
| Elapsed Time: 0:00:00 ETA:   0:00:00
    \end{Verbatim}

    \begin{Verbatim}[commandchars=\\\{\}]
 \def\tcRGB{\textcolor[RGB]}\expandafter\tcRGB\expandafter{\detokenize{255,156,0}}{ 36\%} \def\tcRGB{\textcolor[RGB]}\expandafter\tcRGB\expandafter{\detokenize{255,156,0}}{(4 of 11)} |\#\#\#\#\#\#\#\#\#
| Elapsed Time: 0:00:00 ETA:   0:00:00
    \end{Verbatim}

    \begin{Verbatim}[commandchars=\\\{\}]
 \def\tcRGB{\textcolor[RGB]}\expandafter\tcRGB\expandafter{\detokenize{255,189,0}}{ 45\%} \def\tcRGB{\textcolor[RGB]}\expandafter\tcRGB\expandafter{\detokenize{255,189,0}}{(5 of 11)} |\#\#\#\#\#\#\#\#\#\#\#
| Elapsed Time: 0:00:00 ETA:   0:00:00
    \end{Verbatim}

    \begin{Verbatim}[commandchars=\\\{\}]
 \def\tcRGB{\textcolor[RGB]}\expandafter\tcRGB\expandafter{\detokenize{236,255,0}}{ 63\%} \def\tcRGB{\textcolor[RGB]}\expandafter\tcRGB\expandafter{\detokenize{236,255,0}}{(7 of 11)} |\#\#\#\#\#\#\#\#\#\#\#\#\#\#\#
| Elapsed Time: 0:00:00 ETA:   0:00:00
    \end{Verbatim}

    \begin{Verbatim}[commandchars=\\\{\}]
 \def\tcRGB{\textcolor[RGB]}\expandafter\tcRGB\expandafter{\detokenize{185,255,0}}{ 81\%} \def\tcRGB{\textcolor[RGB]}\expandafter\tcRGB\expandafter{\detokenize{185,255,0}}{(9 of 11)}
|\#\#\#\#\#\#\#\#\#\#\#\#\#\#\#\#\#\#\#\#     | Elapsed Time: 0:00:00 ETA:   0:00:00
    \end{Verbatim}

    \begin{Verbatim}[commandchars=\\\{\}]
 \def\tcRGB{\textcolor[RGB]}\expandafter\tcRGB\expandafter{\detokenize{160,255,0}}{ 90\%} \def\tcRGB{\textcolor[RGB]}\expandafter\tcRGB\expandafter{\detokenize{160,255,0}}{(10 of 11)}
|\#\#\#\#\#\#\#\#\#\#\#\#\#\#\#\#\#\#\#\#\#   | Elapsed Time: 0:00:00 ETA:   0:00:00
    \end{Verbatim}

    \begin{Verbatim}[commandchars=\\\{\}]
 \def\tcRGB{\textcolor[RGB]}\expandafter\tcRGB\expandafter{\detokenize{0,255,0}}{100\%} \def\tcRGB{\textcolor[RGB]}\expandafter\tcRGB\expandafter{\detokenize{0,255,0}}{(11 of 11)}
|\#\#\#\#\#\#\#\#\#\#\#\#\#\#\#\#\#\#\#\#\#\#\#\#| Elapsed Time: 0:00:00 Time:  0:00:00
    \end{Verbatim}

    \begin{Verbatim}[commandchars=\\\{\}]

    \end{Verbatim}
 
            
\prompt{Out}{outcolor}{25}{}
    
    \begin{center}
    \adjustimage{max size={0.9\linewidth}{0.9\paperheight}}{artifacts/latex/Ch07-nonlin-lab_files/artifacts/latex/Ch07-nonlin-lab_54_9.png}
    \end{center}
    { \hspace*{\fill} \\}
    

    Alternatively, we can fix the degrees of freedom of the smoothing spline
using a function included in the \texttt{ISLP.pygam} package. Below we
find a value of \(\lambda\) that gives us roughly four degrees of
freedom. We note here that these degrees of freedom include the
unpenalized intercept and linear term of the smoothing spline, hence
there are at least two degrees of freedom.

    \begin{tcolorbox}[breakable, size=fbox, boxrule=1pt, pad at break*=1mm,colback=cellbackground, colframe=cellborder]
\prompt{In}{incolor}{26}{\boxspacing}
\begin{Verbatim}[commandchars=\\\{\}]
\PY{n}{age\PYZus{}term} \PY{o}{=} \PY{n}{gam}\PY{o}{.}\PY{n}{terms}\PY{p}{[}\PY{l+m+mi}{0}\PY{p}{]}
\PY{n}{lam\PYZus{}4} \PY{o}{=} \PY{n}{approx\PYZus{}lam}\PY{p}{(}\PY{n}{X\PYZus{}age}\PY{p}{,} \PY{n}{age\PYZus{}term}\PY{p}{,} \PY{l+m+mi}{4}\PY{p}{)}
\PY{n}{age\PYZus{}term}\PY{o}{.}\PY{n}{lam} \PY{o}{=} \PY{n}{lam\PYZus{}4}
\PY{n}{degrees\PYZus{}of\PYZus{}freedom}\PY{p}{(}\PY{n}{X\PYZus{}age}\PY{p}{,} \PY{n}{age\PYZus{}term}\PY{p}{)}
\end{Verbatim}
\end{tcolorbox}

            \begin{tcolorbox}[breakable, size=fbox, boxrule=.5pt, pad at break*=1mm, opacityfill=0]
\prompt{Out}{outcolor}{26}{\boxspacing}
\begin{Verbatim}[commandchars=\\\{\}]
4.000000100001664
\end{Verbatim}
\end{tcolorbox}
        
    Let's vary the degrees of freedom in a similar plot to above. We choose
the degrees of freedom as the desired degrees of freedom plus one to
account for the fact that these smoothing splines always have an
intercept term. Hence, a value of one for \texttt{df} is just a linear
fit.

    \begin{tcolorbox}[breakable, size=fbox, boxrule=1pt, pad at break*=1mm,colback=cellbackground, colframe=cellborder]
\prompt{In}{incolor}{27}{\boxspacing}
\begin{Verbatim}[commandchars=\\\{\}]
\PY{n}{fig}\PY{p}{,} \PY{n}{ax} \PY{o}{=} \PY{n}{subplots}\PY{p}{(}\PY{n}{figsize}\PY{o}{=}\PY{p}{(}\PY{l+m+mi}{8}\PY{p}{,}\PY{l+m+mi}{8}\PY{p}{)}\PY{p}{)}
\PY{n}{ax}\PY{o}{.}\PY{n}{scatter}\PY{p}{(}\PY{n}{X\PYZus{}age}\PY{p}{,}
           \PY{n}{y}\PY{p}{,}
           \PY{n}{facecolor}\PY{o}{=}\PY{l+s+s1}{\PYZsq{}}\PY{l+s+s1}{gray}\PY{l+s+s1}{\PYZsq{}}\PY{p}{,}
           \PY{n}{alpha}\PY{o}{=}\PY{l+m+mf}{0.3}\PY{p}{)}
\PY{k}{for} \PY{n}{df} \PY{o+ow}{in} \PY{p}{[}\PY{l+m+mi}{1}\PY{p}{,}\PY{l+m+mi}{3}\PY{p}{,}\PY{l+m+mi}{4}\PY{p}{,}\PY{l+m+mi}{8}\PY{p}{,}\PY{l+m+mi}{15}\PY{p}{]}\PY{p}{:}
    \PY{n}{lam} \PY{o}{=} \PY{n}{approx\PYZus{}lam}\PY{p}{(}\PY{n}{X\PYZus{}age}\PY{p}{,} \PY{n}{age\PYZus{}term}\PY{p}{,} \PY{n}{df}\PY{o}{+}\PY{l+m+mi}{1}\PY{p}{)}
    \PY{n}{age\PYZus{}term}\PY{o}{.}\PY{n}{lam} \PY{o}{=} \PY{n}{lam}
    \PY{n}{gam}\PY{o}{.}\PY{n}{fit}\PY{p}{(}\PY{n}{X\PYZus{}age}\PY{p}{,} \PY{n}{y}\PY{p}{)}
    \PY{n}{ax}\PY{o}{.}\PY{n}{plot}\PY{p}{(}\PY{n}{age\PYZus{}grid}\PY{p}{,}
            \PY{n}{gam}\PY{o}{.}\PY{n}{predict}\PY{p}{(}\PY{n}{age\PYZus{}grid}\PY{p}{)}\PY{p}{,}
            \PY{n}{label}\PY{o}{=}\PY{l+s+s1}{\PYZsq{}}\PY{l+s+si}{\PYZob{}:d\PYZcb{}}\PY{l+s+s1}{\PYZsq{}}\PY{o}{.}\PY{n}{format}\PY{p}{(}\PY{n}{df}\PY{p}{)}\PY{p}{,}
            \PY{n}{linewidth}\PY{o}{=}\PY{l+m+mi}{4}\PY{p}{)}
\PY{n}{ax}\PY{o}{.}\PY{n}{set\PYZus{}xlabel}\PY{p}{(}\PY{l+s+s1}{\PYZsq{}}\PY{l+s+s1}{Age}\PY{l+s+s1}{\PYZsq{}}\PY{p}{,} \PY{n}{fontsize}\PY{o}{=}\PY{l+m+mi}{20}\PY{p}{)}
\PY{n}{ax}\PY{o}{.}\PY{n}{set\PYZus{}ylabel}\PY{p}{(}\PY{l+s+s1}{\PYZsq{}}\PY{l+s+s1}{Wage}\PY{l+s+s1}{\PYZsq{}}\PY{p}{,} \PY{n}{fontsize}\PY{o}{=}\PY{l+m+mi}{20}\PY{p}{)}\PY{p}{;}
\PY{n}{ax}\PY{o}{.}\PY{n}{legend}\PY{p}{(}\PY{n}{title}\PY{o}{=}\PY{l+s+s1}{\PYZsq{}}\PY{l+s+s1}{Degrees of freedom}\PY{l+s+s1}{\PYZsq{}}\PY{p}{)}\PY{p}{;}
\end{Verbatim}
\end{tcolorbox}

    \begin{center}
    \adjustimage{max size={0.9\linewidth}{0.9\paperheight}}{artifacts/latex/Ch07-nonlin-lab_files/artifacts/latex/Ch07-nonlin-lab_58_0.png}
    \end{center}
    { \hspace*{\fill} \\}
    
    \hypertarget{additive-models-with-several-terms}{%
\subsubsection{Additive Models with Several
Terms}\label{additive-models-with-several-terms}}

The strength of generalized additive models lies in their ability to fit
multivariate regression models with more flexibility than linear models.
We demonstrate two approaches: the first in a more manual fashion using
natural splines and piecewise constant functions, and the second using
the \texttt{pygam} package and smoothing splines.

We now fit a GAM by hand to predict \texttt{wage} using natural spline
functions of \texttt{year} and \texttt{age}, treating \texttt{education}
as a qualitative predictor, as in (7.16). Since this is just a big
linear regression model using an appropriate choice of basis functions,
we can simply do this using the \texttt{sm.OLS()} function.

We will build the model matrix in a more manual fashion here, since we
wish to access the pieces separately when constructing partial
dependence plots.

    \begin{tcolorbox}[breakable, size=fbox, boxrule=1pt, pad at break*=1mm,colback=cellbackground, colframe=cellborder]
\prompt{In}{incolor}{28}{\boxspacing}
\begin{Verbatim}[commandchars=\\\{\}]
\PY{n}{ns\PYZus{}age} \PY{o}{=} \PY{n}{NaturalSpline}\PY{p}{(}\PY{n}{df}\PY{o}{=}\PY{l+m+mi}{4}\PY{p}{)}\PY{o}{.}\PY{n}{fit}\PY{p}{(}\PY{n}{age}\PY{p}{)}
\PY{n}{ns\PYZus{}year} \PY{o}{=} \PY{n}{NaturalSpline}\PY{p}{(}\PY{n}{df}\PY{o}{=}\PY{l+m+mi}{5}\PY{p}{)}\PY{o}{.}\PY{n}{fit}\PY{p}{(}\PY{n}{Wage}\PY{p}{[}\PY{l+s+s1}{\PYZsq{}}\PY{l+s+s1}{year}\PY{l+s+s1}{\PYZsq{}}\PY{p}{]}\PY{p}{)}
\PY{n}{Xs} \PY{o}{=} \PY{p}{[}\PY{n}{ns\PYZus{}age}\PY{o}{.}\PY{n}{transform}\PY{p}{(}\PY{n}{age}\PY{p}{)}\PY{p}{,}
      \PY{n}{ns\PYZus{}year}\PY{o}{.}\PY{n}{transform}\PY{p}{(}\PY{n}{Wage}\PY{p}{[}\PY{l+s+s1}{\PYZsq{}}\PY{l+s+s1}{year}\PY{l+s+s1}{\PYZsq{}}\PY{p}{]}\PY{p}{)}\PY{p}{,}
      \PY{n}{pd}\PY{o}{.}\PY{n}{get\PYZus{}dummies}\PY{p}{(}\PY{n}{Wage}\PY{p}{[}\PY{l+s+s1}{\PYZsq{}}\PY{l+s+s1}{education}\PY{l+s+s1}{\PYZsq{}}\PY{p}{]}\PY{p}{)}\PY{o}{.}\PY{n}{values}\PY{p}{]}
\PY{n}{X\PYZus{}bh} \PY{o}{=} \PY{n}{np}\PY{o}{.}\PY{n}{hstack}\PY{p}{(}\PY{n}{Xs}\PY{p}{)}
\PY{n}{gam\PYZus{}bh} \PY{o}{=} \PY{n}{sm}\PY{o}{.}\PY{n}{OLS}\PY{p}{(}\PY{n}{y}\PY{p}{,} \PY{n}{X\PYZus{}bh}\PY{p}{)}\PY{o}{.}\PY{n}{fit}\PY{p}{(}\PY{p}{)}
\end{Verbatim}
\end{tcolorbox}

    Here the function \texttt{NaturalSpline()} is the workhorse supporting
the \texttt{ns()} helper function. We chose to use all columns of the
indicator matrix for the categorical variable \texttt{education}, making
an intercept redundant. Finally, we stacked the three component matrices
horizontally to form the model matrix \texttt{X\_bh}.

We now show how to construct partial dependence plots for each of the
terms in our rudimentary GAM. We can do this by hand, given grids for
\texttt{age} and \texttt{year}. We simply predict with new \(X\)
matrices, fixing all but one of the features at a time.

    \begin{tcolorbox}[breakable, size=fbox, boxrule=1pt, pad at break*=1mm,colback=cellbackground, colframe=cellborder]
\prompt{In}{incolor}{29}{\boxspacing}
\begin{Verbatim}[commandchars=\\\{\}]
\PY{n}{age\PYZus{}grid} \PY{o}{=} \PY{n}{np}\PY{o}{.}\PY{n}{linspace}\PY{p}{(}\PY{n}{age}\PY{o}{.}\PY{n}{min}\PY{p}{(}\PY{p}{)}\PY{p}{,}
                       \PY{n}{age}\PY{o}{.}\PY{n}{max}\PY{p}{(}\PY{p}{)}\PY{p}{,}
                       \PY{l+m+mi}{100}\PY{p}{)}
\PY{n}{X\PYZus{}age\PYZus{}bh} \PY{o}{=} \PY{n}{X\PYZus{}bh}\PY{o}{.}\PY{n}{copy}\PY{p}{(}\PY{p}{)}\PY{p}{[}\PY{p}{:}\PY{l+m+mi}{100}\PY{p}{]}
\PY{n}{X\PYZus{}age\PYZus{}bh}\PY{p}{[}\PY{p}{:}\PY{p}{]} \PY{o}{=} \PY{n}{X\PYZus{}bh}\PY{p}{[}\PY{p}{:}\PY{p}{]}\PY{o}{.}\PY{n}{mean}\PY{p}{(}\PY{l+m+mi}{0}\PY{p}{)}\PY{p}{[}\PY{k+kc}{None}\PY{p}{,}\PY{p}{:}\PY{p}{]}
\PY{n}{X\PYZus{}age\PYZus{}bh}\PY{p}{[}\PY{p}{:}\PY{p}{,}\PY{p}{:}\PY{l+m+mi}{4}\PY{p}{]} \PY{o}{=} \PY{n}{ns\PYZus{}age}\PY{o}{.}\PY{n}{transform}\PY{p}{(}\PY{n}{age\PYZus{}grid}\PY{p}{)}
\PY{n}{preds} \PY{o}{=} \PY{n}{gam\PYZus{}bh}\PY{o}{.}\PY{n}{get\PYZus{}prediction}\PY{p}{(}\PY{n}{X\PYZus{}age\PYZus{}bh}\PY{p}{)}
\PY{n}{bounds\PYZus{}age} \PY{o}{=} \PY{n}{preds}\PY{o}{.}\PY{n}{conf\PYZus{}int}\PY{p}{(}\PY{n}{alpha}\PY{o}{=}\PY{l+m+mf}{0.05}\PY{p}{)}
\PY{n}{partial\PYZus{}age} \PY{o}{=} \PY{n}{preds}\PY{o}{.}\PY{n}{predicted\PYZus{}mean}
\PY{n}{center} \PY{o}{=} \PY{n}{partial\PYZus{}age}\PY{o}{.}\PY{n}{mean}\PY{p}{(}\PY{p}{)}
\PY{n}{partial\PYZus{}age} \PY{o}{\PYZhy{}}\PY{o}{=} \PY{n}{center}
\PY{n}{bounds\PYZus{}age} \PY{o}{\PYZhy{}}\PY{o}{=} \PY{n}{center}
\PY{n}{fig}\PY{p}{,} \PY{n}{ax} \PY{o}{=} \PY{n}{subplots}\PY{p}{(}\PY{n}{figsize}\PY{o}{=}\PY{p}{(}\PY{l+m+mi}{8}\PY{p}{,}\PY{l+m+mi}{8}\PY{p}{)}\PY{p}{)}
\PY{n}{ax}\PY{o}{.}\PY{n}{plot}\PY{p}{(}\PY{n}{age\PYZus{}grid}\PY{p}{,} \PY{n}{partial\PYZus{}age}\PY{p}{,} \PY{l+s+s1}{\PYZsq{}}\PY{l+s+s1}{b}\PY{l+s+s1}{\PYZsq{}}\PY{p}{,} \PY{n}{linewidth}\PY{o}{=}\PY{l+m+mi}{3}\PY{p}{)}
\PY{n}{ax}\PY{o}{.}\PY{n}{plot}\PY{p}{(}\PY{n}{age\PYZus{}grid}\PY{p}{,} \PY{n}{bounds\PYZus{}age}\PY{p}{[}\PY{p}{:}\PY{p}{,}\PY{l+m+mi}{0}\PY{p}{]}\PY{p}{,} \PY{l+s+s1}{\PYZsq{}}\PY{l+s+s1}{r\PYZhy{}\PYZhy{}}\PY{l+s+s1}{\PYZsq{}}\PY{p}{,} \PY{n}{linewidth}\PY{o}{=}\PY{l+m+mi}{3}\PY{p}{)}
\PY{n}{ax}\PY{o}{.}\PY{n}{plot}\PY{p}{(}\PY{n}{age\PYZus{}grid}\PY{p}{,} \PY{n}{bounds\PYZus{}age}\PY{p}{[}\PY{p}{:}\PY{p}{,}\PY{l+m+mi}{1}\PY{p}{]}\PY{p}{,} \PY{l+s+s1}{\PYZsq{}}\PY{l+s+s1}{r\PYZhy{}\PYZhy{}}\PY{l+s+s1}{\PYZsq{}}\PY{p}{,} \PY{n}{linewidth}\PY{o}{=}\PY{l+m+mi}{3}\PY{p}{)}
\PY{n}{ax}\PY{o}{.}\PY{n}{set\PYZus{}xlabel}\PY{p}{(}\PY{l+s+s1}{\PYZsq{}}\PY{l+s+s1}{Age}\PY{l+s+s1}{\PYZsq{}}\PY{p}{)}
\PY{n}{ax}\PY{o}{.}\PY{n}{set\PYZus{}ylabel}\PY{p}{(}\PY{l+s+s1}{\PYZsq{}}\PY{l+s+s1}{Effect on wage}\PY{l+s+s1}{\PYZsq{}}\PY{p}{)}
\PY{n}{ax}\PY{o}{.}\PY{n}{set\PYZus{}title}\PY{p}{(}\PY{l+s+s1}{\PYZsq{}}\PY{l+s+s1}{Partial dependence of age on wage}\PY{l+s+s1}{\PYZsq{}}\PY{p}{,} \PY{n}{fontsize}\PY{o}{=}\PY{l+m+mi}{20}\PY{p}{)}\PY{p}{;}
\end{Verbatim}
\end{tcolorbox}

    \begin{center}
    \adjustimage{max size={0.9\linewidth}{0.9\paperheight}}{artifacts/latex/Ch07-nonlin-lab_files/artifacts/latex/Ch07-nonlin-lab_62_0.png}
    \end{center}
    { \hspace*{\fill} \\}
    
    Let's explain in some detail what we did above. The idea is to create a
new prediction matrix, where all but the columns belonging to
\texttt{age} are constant (and set to their training-data means). The
four columns for \texttt{age} are filled in with the natural spline
basis evaluated at the 100 values in \texttt{age\_grid}.

\begin{itemize}
\tightlist
\item
  We made a grid of length 100 in \texttt{age}, and created a matrix
  \texttt{X\_age\_bh} with 100 rows and the same number of columns as
  \texttt{X\_bh}.
\item
  We replaced every row of this matrix with the column means of the
  original.
\item
  We then replace just the first four columns representing \texttt{age}
  with the natural spline basis computed at the values in
  \texttt{age\_grid}.
\end{itemize}

The remaining steps should by now be familiar.

We also look at the effect of \texttt{year} on \texttt{wage}; the
process is the same.

    \begin{tcolorbox}[breakable, size=fbox, boxrule=1pt, pad at break*=1mm,colback=cellbackground, colframe=cellborder]
\prompt{In}{incolor}{30}{\boxspacing}
\begin{Verbatim}[commandchars=\\\{\}]
\PY{n}{year\PYZus{}grid} \PY{o}{=} \PY{n}{np}\PY{o}{.}\PY{n}{linspace}\PY{p}{(}\PY{l+m+mi}{2003}\PY{p}{,} \PY{l+m+mi}{2009}\PY{p}{,} \PY{l+m+mi}{100}\PY{p}{)}
\PY{n}{year\PYZus{}grid} \PY{o}{=} \PY{n}{np}\PY{o}{.}\PY{n}{linspace}\PY{p}{(}\PY{n}{Wage}\PY{p}{[}\PY{l+s+s1}{\PYZsq{}}\PY{l+s+s1}{year}\PY{l+s+s1}{\PYZsq{}}\PY{p}{]}\PY{o}{.}\PY{n}{min}\PY{p}{(}\PY{p}{)}\PY{p}{,}
                        \PY{n}{Wage}\PY{p}{[}\PY{l+s+s1}{\PYZsq{}}\PY{l+s+s1}{year}\PY{l+s+s1}{\PYZsq{}}\PY{p}{]}\PY{o}{.}\PY{n}{max}\PY{p}{(}\PY{p}{)}\PY{p}{,}
                        \PY{l+m+mi}{100}\PY{p}{)}
\PY{n}{X\PYZus{}year\PYZus{}bh} \PY{o}{=} \PY{n}{X\PYZus{}bh}\PY{o}{.}\PY{n}{copy}\PY{p}{(}\PY{p}{)}\PY{p}{[}\PY{p}{:}\PY{l+m+mi}{100}\PY{p}{]}
\PY{n}{X\PYZus{}year\PYZus{}bh}\PY{p}{[}\PY{p}{:}\PY{p}{]} \PY{o}{=} \PY{n}{X\PYZus{}bh}\PY{p}{[}\PY{p}{:}\PY{p}{]}\PY{o}{.}\PY{n}{mean}\PY{p}{(}\PY{l+m+mi}{0}\PY{p}{)}\PY{p}{[}\PY{k+kc}{None}\PY{p}{,}\PY{p}{:}\PY{p}{]}
\PY{n}{X\PYZus{}year\PYZus{}bh}\PY{p}{[}\PY{p}{:}\PY{p}{,}\PY{l+m+mi}{4}\PY{p}{:}\PY{l+m+mi}{9}\PY{p}{]} \PY{o}{=} \PY{n}{ns\PYZus{}year}\PY{o}{.}\PY{n}{transform}\PY{p}{(}\PY{n}{year\PYZus{}grid}\PY{p}{)}
\PY{n}{preds} \PY{o}{=} \PY{n}{gam\PYZus{}bh}\PY{o}{.}\PY{n}{get\PYZus{}prediction}\PY{p}{(}\PY{n}{X\PYZus{}year\PYZus{}bh}\PY{p}{)}
\PY{n}{bounds\PYZus{}year} \PY{o}{=} \PY{n}{preds}\PY{o}{.}\PY{n}{conf\PYZus{}int}\PY{p}{(}\PY{n}{alpha}\PY{o}{=}\PY{l+m+mf}{0.05}\PY{p}{)}
\PY{n}{partial\PYZus{}year} \PY{o}{=} \PY{n}{preds}\PY{o}{.}\PY{n}{predicted\PYZus{}mean}
\PY{n}{center} \PY{o}{=} \PY{n}{partial\PYZus{}year}\PY{o}{.}\PY{n}{mean}\PY{p}{(}\PY{p}{)}
\PY{n}{partial\PYZus{}year} \PY{o}{\PYZhy{}}\PY{o}{=} \PY{n}{center}
\PY{n}{bounds\PYZus{}year} \PY{o}{\PYZhy{}}\PY{o}{=} \PY{n}{center}
\PY{n}{fig}\PY{p}{,} \PY{n}{ax} \PY{o}{=} \PY{n}{subplots}\PY{p}{(}\PY{n}{figsize}\PY{o}{=}\PY{p}{(}\PY{l+m+mi}{8}\PY{p}{,}\PY{l+m+mi}{8}\PY{p}{)}\PY{p}{)}
\PY{n}{ax}\PY{o}{.}\PY{n}{plot}\PY{p}{(}\PY{n}{year\PYZus{}grid}\PY{p}{,} \PY{n}{partial\PYZus{}year}\PY{p}{,} \PY{l+s+s1}{\PYZsq{}}\PY{l+s+s1}{b}\PY{l+s+s1}{\PYZsq{}}\PY{p}{,} \PY{n}{linewidth}\PY{o}{=}\PY{l+m+mi}{3}\PY{p}{)}
\PY{n}{ax}\PY{o}{.}\PY{n}{plot}\PY{p}{(}\PY{n}{year\PYZus{}grid}\PY{p}{,} \PY{n}{bounds\PYZus{}year}\PY{p}{[}\PY{p}{:}\PY{p}{,}\PY{l+m+mi}{0}\PY{p}{]}\PY{p}{,} \PY{l+s+s1}{\PYZsq{}}\PY{l+s+s1}{r\PYZhy{}\PYZhy{}}\PY{l+s+s1}{\PYZsq{}}\PY{p}{,} \PY{n}{linewidth}\PY{o}{=}\PY{l+m+mi}{3}\PY{p}{)}
\PY{n}{ax}\PY{o}{.}\PY{n}{plot}\PY{p}{(}\PY{n}{year\PYZus{}grid}\PY{p}{,} \PY{n}{bounds\PYZus{}year}\PY{p}{[}\PY{p}{:}\PY{p}{,}\PY{l+m+mi}{1}\PY{p}{]}\PY{p}{,} \PY{l+s+s1}{\PYZsq{}}\PY{l+s+s1}{r\PYZhy{}\PYZhy{}}\PY{l+s+s1}{\PYZsq{}}\PY{p}{,} \PY{n}{linewidth}\PY{o}{=}\PY{l+m+mi}{3}\PY{p}{)}
\PY{n}{ax}\PY{o}{.}\PY{n}{set\PYZus{}xlabel}\PY{p}{(}\PY{l+s+s1}{\PYZsq{}}\PY{l+s+s1}{Year}\PY{l+s+s1}{\PYZsq{}}\PY{p}{)}
\PY{n}{ax}\PY{o}{.}\PY{n}{set\PYZus{}ylabel}\PY{p}{(}\PY{l+s+s1}{\PYZsq{}}\PY{l+s+s1}{Effect on wage}\PY{l+s+s1}{\PYZsq{}}\PY{p}{)}
\PY{n}{ax}\PY{o}{.}\PY{n}{set\PYZus{}title}\PY{p}{(}\PY{l+s+s1}{\PYZsq{}}\PY{l+s+s1}{Partial dependence of year on wage}\PY{l+s+s1}{\PYZsq{}}\PY{p}{,} \PY{n}{fontsize}\PY{o}{=}\PY{l+m+mi}{20}\PY{p}{)}\PY{p}{;}
\end{Verbatim}
\end{tcolorbox}

    \begin{center}
    \adjustimage{max size={0.9\linewidth}{0.9\paperheight}}{artifacts/latex/Ch07-nonlin-lab_files/artifacts/latex/Ch07-nonlin-lab_64_0.png}
    \end{center}
    { \hspace*{\fill} \\}
    
    We now fit the model (7.16) using smoothing splines rather than natural
splines. All of the terms in (7.16) are fit simultaneously, taking each
other into account to explain the response. The \texttt{pygam} package
only works with matrices, so we must convert the categorical series
\texttt{education} to its array representation, which can be found with
the \texttt{cat.codes} attribute of \texttt{education}. As \texttt{year}
only has 7 unique values, we use only seven basis functions for it.

    \begin{tcolorbox}[breakable, size=fbox, boxrule=1pt, pad at break*=1mm,colback=cellbackground, colframe=cellborder]
\prompt{In}{incolor}{31}{\boxspacing}
\begin{Verbatim}[commandchars=\\\{\}]
\PY{n}{gam\PYZus{}full} \PY{o}{=} \PY{n}{LinearGAM}\PY{p}{(}\PY{n}{s\PYZus{}gam}\PY{p}{(}\PY{l+m+mi}{0}\PY{p}{)} \PY{o}{+}
                     \PY{n}{s\PYZus{}gam}\PY{p}{(}\PY{l+m+mi}{1}\PY{p}{,} \PY{n}{n\PYZus{}splines}\PY{o}{=}\PY{l+m+mi}{7}\PY{p}{)} \PY{o}{+}
                     \PY{n}{f\PYZus{}gam}\PY{p}{(}\PY{l+m+mi}{2}\PY{p}{,} \PY{n}{lam}\PY{o}{=}\PY{l+m+mi}{0}\PY{p}{)}\PY{p}{)}
\PY{n}{Xgam} \PY{o}{=} \PY{n}{np}\PY{o}{.}\PY{n}{column\PYZus{}stack}\PY{p}{(}\PY{p}{[}\PY{n}{age}\PY{p}{,}
                        \PY{n}{Wage}\PY{p}{[}\PY{l+s+s1}{\PYZsq{}}\PY{l+s+s1}{year}\PY{l+s+s1}{\PYZsq{}}\PY{p}{]}\PY{p}{,}
                        \PY{n}{Wage}\PY{p}{[}\PY{l+s+s1}{\PYZsq{}}\PY{l+s+s1}{education}\PY{l+s+s1}{\PYZsq{}}\PY{p}{]}\PY{o}{.}\PY{n}{cat}\PY{o}{.}\PY{n}{codes}\PY{p}{]}\PY{p}{)}
\PY{n}{gam\PYZus{}full} \PY{o}{=} \PY{n}{gam\PYZus{}full}\PY{o}{.}\PY{n}{fit}\PY{p}{(}\PY{n}{Xgam}\PY{p}{,} \PY{n}{y}\PY{p}{)}
\end{Verbatim}
\end{tcolorbox}

    The two \texttt{s\_gam()} terms result in smoothing spline fits, and use
a default value for \(\lambda\) (\texttt{lam=0.6}), which is somewhat
arbitrary. For the categorical term \texttt{education}, specified using
a \texttt{f\_gam()} term, we specify \texttt{lam=0} to avoid any
shrinkage. We produce the partial dependence plot in \texttt{age} to see
the effect of these choices.

The values for the plot are generated by the \texttt{pygam} package. We
provide a \texttt{plot\_gam()} function for partial-dependence plots in
\texttt{ISLP.pygam}, which makes this job easier than in our last
example with natural splines.

    \begin{tcolorbox}[breakable, size=fbox, boxrule=1pt, pad at break*=1mm,colback=cellbackground, colframe=cellborder]
\prompt{In}{incolor}{32}{\boxspacing}
\begin{Verbatim}[commandchars=\\\{\}]
\PY{n}{fig}\PY{p}{,} \PY{n}{ax} \PY{o}{=} \PY{n}{subplots}\PY{p}{(}\PY{n}{figsize}\PY{o}{=}\PY{p}{(}\PY{l+m+mi}{8}\PY{p}{,}\PY{l+m+mi}{8}\PY{p}{)}\PY{p}{)}
\PY{n}{plot\PYZus{}gam}\PY{p}{(}\PY{n}{gam\PYZus{}full}\PY{p}{,} \PY{l+m+mi}{0}\PY{p}{,} \PY{n}{ax}\PY{o}{=}\PY{n}{ax}\PY{p}{)}
\PY{n}{ax}\PY{o}{.}\PY{n}{set\PYZus{}xlabel}\PY{p}{(}\PY{l+s+s1}{\PYZsq{}}\PY{l+s+s1}{Age}\PY{l+s+s1}{\PYZsq{}}\PY{p}{)}
\PY{n}{ax}\PY{o}{.}\PY{n}{set\PYZus{}ylabel}\PY{p}{(}\PY{l+s+s1}{\PYZsq{}}\PY{l+s+s1}{Effect on wage}\PY{l+s+s1}{\PYZsq{}}\PY{p}{)}
\PY{n}{ax}\PY{o}{.}\PY{n}{set\PYZus{}title}\PY{p}{(}\PY{l+s+s1}{\PYZsq{}}\PY{l+s+s1}{Partial dependence of age on wage \PYZhy{} default lam=0.6}\PY{l+s+s1}{\PYZsq{}}\PY{p}{,} \PY{n}{fontsize}\PY{o}{=}\PY{l+m+mi}{20}\PY{p}{)}\PY{p}{;}
\end{Verbatim}
\end{tcolorbox}

    \begin{center}
    \adjustimage{max size={0.9\linewidth}{0.9\paperheight}}{artifacts/latex/Ch07-nonlin-lab_files/artifacts/latex/Ch07-nonlin-lab_68_0.png}
    \end{center}
    { \hspace*{\fill} \\}
    
    We see that the function is somewhat wiggly. It is more natural to
specify the \texttt{df} than a value for \texttt{lam}. We refit a GAM
using four degrees of freedom each for \texttt{age} and \texttt{year}.
Recall that the addition of one below takes into account the intercept
of the smoothing spline.

    \begin{tcolorbox}[breakable, size=fbox, boxrule=1pt, pad at break*=1mm,colback=cellbackground, colframe=cellborder]
\prompt{In}{incolor}{33}{\boxspacing}
\begin{Verbatim}[commandchars=\\\{\}]
\PY{n}{age\PYZus{}term} \PY{o}{=} \PY{n}{gam\PYZus{}full}\PY{o}{.}\PY{n}{terms}\PY{p}{[}\PY{l+m+mi}{0}\PY{p}{]}
\PY{n}{age\PYZus{}term}\PY{o}{.}\PY{n}{lam} \PY{o}{=} \PY{n}{approx\PYZus{}lam}\PY{p}{(}\PY{n}{Xgam}\PY{p}{,} \PY{n}{age\PYZus{}term}\PY{p}{,} \PY{n}{df}\PY{o}{=}\PY{l+m+mi}{4}\PY{o}{+}\PY{l+m+mi}{1}\PY{p}{)}
\PY{n}{year\PYZus{}term} \PY{o}{=} \PY{n}{gam\PYZus{}full}\PY{o}{.}\PY{n}{terms}\PY{p}{[}\PY{l+m+mi}{1}\PY{p}{]}
\PY{n}{year\PYZus{}term}\PY{o}{.}\PY{n}{lam} \PY{o}{=} \PY{n}{approx\PYZus{}lam}\PY{p}{(}\PY{n}{Xgam}\PY{p}{,} \PY{n}{year\PYZus{}term}\PY{p}{,} \PY{n}{df}\PY{o}{=}\PY{l+m+mi}{4}\PY{o}{+}\PY{l+m+mi}{1}\PY{p}{)}
\PY{n}{gam\PYZus{}full} \PY{o}{=} \PY{n}{gam\PYZus{}full}\PY{o}{.}\PY{n}{fit}\PY{p}{(}\PY{n}{Xgam}\PY{p}{,} \PY{n}{y}\PY{p}{)}
\end{Verbatim}
\end{tcolorbox}

    Note that updating \texttt{age\_term.lam} above updates it in
\texttt{gam\_full.terms{[}0{]}} as well! Likewise for
\texttt{year\_term.lam}.

Repeating the plot for \texttt{age}, we see that it is much smoother. We
also produce the plot for \texttt{year}.

    \begin{tcolorbox}[breakable, size=fbox, boxrule=1pt, pad at break*=1mm,colback=cellbackground, colframe=cellborder]
\prompt{In}{incolor}{34}{\boxspacing}
\begin{Verbatim}[commandchars=\\\{\}]
\PY{n}{fig}\PY{p}{,} \PY{n}{ax} \PY{o}{=} \PY{n}{subplots}\PY{p}{(}\PY{n}{figsize}\PY{o}{=}\PY{p}{(}\PY{l+m+mi}{8}\PY{p}{,}\PY{l+m+mi}{8}\PY{p}{)}\PY{p}{)}
\PY{n}{plot\PYZus{}gam}\PY{p}{(}\PY{n}{gam\PYZus{}full}\PY{p}{,}
         \PY{l+m+mi}{1}\PY{p}{,}
         \PY{n}{ax}\PY{o}{=}\PY{n}{ax}\PY{p}{)}
\PY{n}{ax}\PY{o}{.}\PY{n}{set\PYZus{}xlabel}\PY{p}{(}\PY{l+s+s1}{\PYZsq{}}\PY{l+s+s1}{Year}\PY{l+s+s1}{\PYZsq{}}\PY{p}{)}
\PY{n}{ax}\PY{o}{.}\PY{n}{set\PYZus{}ylabel}\PY{p}{(}\PY{l+s+s1}{\PYZsq{}}\PY{l+s+s1}{Effect on wage}\PY{l+s+s1}{\PYZsq{}}\PY{p}{)}
\PY{n}{ax}\PY{o}{.}\PY{n}{set\PYZus{}title}\PY{p}{(}\PY{l+s+s1}{\PYZsq{}}\PY{l+s+s1}{Partial dependence of year on wage}\PY{l+s+s1}{\PYZsq{}}\PY{p}{,} \PY{n}{fontsize}\PY{o}{=}\PY{l+m+mi}{20}\PY{p}{)}
\end{Verbatim}
\end{tcolorbox}

            \begin{tcolorbox}[breakable, size=fbox, boxrule=.5pt, pad at break*=1mm, opacityfill=0]
\prompt{Out}{outcolor}{34}{\boxspacing}
\begin{Verbatim}[commandchars=\\\{\}]
Text(0.5, 1.0, 'Partial dependence of year on wage')
\end{Verbatim}
\end{tcolorbox}
        
    \begin{center}
    \adjustimage{max size={0.9\linewidth}{0.9\paperheight}}{artifacts/latex/Ch07-nonlin-lab_files/artifacts/latex/Ch07-nonlin-lab_72_1.png}
    \end{center}
    { \hspace*{\fill} \\}
    
    Finally we plot \texttt{education}, which is categorical. The partial
dependence plot is different, and more suitable for the set of fitted
constants for each level of this variable.

    \begin{tcolorbox}[breakable, size=fbox, boxrule=1pt, pad at break*=1mm,colback=cellbackground, colframe=cellborder]
\prompt{In}{incolor}{35}{\boxspacing}
\begin{Verbatim}[commandchars=\\\{\}]
\PY{n}{fig}\PY{p}{,} \PY{n}{ax} \PY{o}{=} \PY{n}{subplots}\PY{p}{(}\PY{n}{figsize}\PY{o}{=}\PY{p}{(}\PY{l+m+mi}{8}\PY{p}{,} \PY{l+m+mi}{8}\PY{p}{)}\PY{p}{)}
\PY{n}{ax} \PY{o}{=} \PY{n}{plot\PYZus{}gam}\PY{p}{(}\PY{n}{gam\PYZus{}full}\PY{p}{,} \PY{l+m+mi}{2}\PY{p}{)}
\PY{n}{ax}\PY{o}{.}\PY{n}{set\PYZus{}xlabel}\PY{p}{(}\PY{l+s+s1}{\PYZsq{}}\PY{l+s+s1}{Education}\PY{l+s+s1}{\PYZsq{}}\PY{p}{)}
\PY{n}{ax}\PY{o}{.}\PY{n}{set\PYZus{}ylabel}\PY{p}{(}\PY{l+s+s1}{\PYZsq{}}\PY{l+s+s1}{Effect on wage}\PY{l+s+s1}{\PYZsq{}}\PY{p}{)}
\PY{n}{ax}\PY{o}{.}\PY{n}{set\PYZus{}title}\PY{p}{(}\PY{l+s+s1}{\PYZsq{}}\PY{l+s+s1}{Partial dependence of wage on education}\PY{l+s+s1}{\PYZsq{}}\PY{p}{,}
             \PY{n}{fontsize}\PY{o}{=}\PY{l+m+mi}{20}\PY{p}{)}\PY{p}{;}
\PY{n}{ax}\PY{o}{.}\PY{n}{set\PYZus{}xticklabels}\PY{p}{(}\PY{n}{Wage}\PY{p}{[}\PY{l+s+s1}{\PYZsq{}}\PY{l+s+s1}{education}\PY{l+s+s1}{\PYZsq{}}\PY{p}{]}\PY{o}{.}\PY{n}{cat}\PY{o}{.}\PY{n}{categories}\PY{p}{,} \PY{n}{fontsize}\PY{o}{=}\PY{l+m+mi}{8}\PY{p}{)}\PY{p}{;}
\end{Verbatim}
\end{tcolorbox}

    \begin{center}
    \adjustimage{max size={0.9\linewidth}{0.9\paperheight}}{artifacts/latex/Ch07-nonlin-lab_files/artifacts/latex/Ch07-nonlin-lab_74_0.png}
    \end{center}
    { \hspace*{\fill} \\}
    
    \hypertarget{anova-tests-for-additive-models}{%
\subsubsection{ANOVA Tests for Additive
Models}\label{anova-tests-for-additive-models}}

In all of our models, the function of \texttt{year} looks rather linear.
We can perform a series of ANOVA tests in order to determine which of
these three models is best: a GAM that excludes \texttt{year}
(\(\mathcal{M}_1\)), a GAM that uses a linear function of \texttt{year}
(\(\mathcal{M}_2\)), or a GAM that uses a spline function of
\texttt{year} (\(\mathcal{M}_3\)).

    \begin{tcolorbox}[breakable, size=fbox, boxrule=1pt, pad at break*=1mm,colback=cellbackground, colframe=cellborder]
\prompt{In}{incolor}{36}{\boxspacing}
\begin{Verbatim}[commandchars=\\\{\}]
\PY{n}{gam\PYZus{}0} \PY{o}{=} \PY{n}{LinearGAM}\PY{p}{(}\PY{n}{age\PYZus{}term} \PY{o}{+} \PY{n}{f\PYZus{}gam}\PY{p}{(}\PY{l+m+mi}{2}\PY{p}{,} \PY{n}{lam}\PY{o}{=}\PY{l+m+mi}{0}\PY{p}{)}\PY{p}{)}
\PY{n}{gam\PYZus{}0}\PY{o}{.}\PY{n}{fit}\PY{p}{(}\PY{n}{Xgam}\PY{p}{,} \PY{n}{y}\PY{p}{)}
\PY{n}{gam\PYZus{}linear} \PY{o}{=} \PY{n}{LinearGAM}\PY{p}{(}\PY{n}{age\PYZus{}term} \PY{o}{+}
                       \PY{n}{l\PYZus{}gam}\PY{p}{(}\PY{l+m+mi}{1}\PY{p}{,} \PY{n}{lam}\PY{o}{=}\PY{l+m+mi}{0}\PY{p}{)} \PY{o}{+}
                       \PY{n}{f\PYZus{}gam}\PY{p}{(}\PY{l+m+mi}{2}\PY{p}{,} \PY{n}{lam}\PY{o}{=}\PY{l+m+mi}{0}\PY{p}{)}\PY{p}{)}
\PY{n}{gam\PYZus{}linear}\PY{o}{.}\PY{n}{fit}\PY{p}{(}\PY{n}{Xgam}\PY{p}{,} \PY{n}{y}\PY{p}{)}
\end{Verbatim}
\end{tcolorbox}

            \begin{tcolorbox}[breakable, size=fbox, boxrule=.5pt, pad at break*=1mm, opacityfill=0]
\prompt{Out}{outcolor}{36}{\boxspacing}
\begin{Verbatim}[commandchars=\\\{\}]
LinearGAM(callbacks=[Deviance(), Diffs()], fit\_intercept=True,
   max\_iter=100, scale=None, terms=s(0) + l(1) + f(2) + intercept,
   tol=0.0001, verbose=False)
\end{Verbatim}
\end{tcolorbox}
        
    Notice our use of \texttt{age\_term} in the expressions above. We do
this because earlier we set the value for \texttt{lam} in this term to
achieve four degrees of freedom.

To directly assess the effect of \texttt{year} we run an ANOVA on the
three models fit above.

    \begin{tcolorbox}[breakable, size=fbox, boxrule=1pt, pad at break*=1mm,colback=cellbackground, colframe=cellborder]
\prompt{In}{incolor}{37}{\boxspacing}
\begin{Verbatim}[commandchars=\\\{\}]
\PY{n}{anova\PYZus{}gam}\PY{p}{(}\PY{n}{gam\PYZus{}0}\PY{p}{,} \PY{n}{gam\PYZus{}linear}\PY{p}{,} \PY{n}{gam\PYZus{}full}\PY{p}{)}
\end{Verbatim}
\end{tcolorbox}

            \begin{tcolorbox}[breakable, size=fbox, boxrule=.5pt, pad at break*=1mm, opacityfill=0]
\prompt{Out}{outcolor}{37}{\boxspacing}
\begin{Verbatim}[commandchars=\\\{\}]
       deviance           df  deviance\_diff   df\_diff          F    pvalue
0  3.714362e+06  2991.004005            NaN       NaN        NaN       NaN
1  3.696746e+06  2990.005190   17616.542840  0.998815  14.265131  0.002314
2  3.693143e+06  2987.007254    3602.893655  2.997936   0.972007  0.435579
\end{Verbatim}
\end{tcolorbox}
        
    We find that there is compelling evidence that a GAM with a linear
function in \texttt{year} is better than a GAM that does not include
\texttt{year} at all (\(p\)-value= 0.002). However, there is no evidence
that a non-linear function of \texttt{year} is needed
(\(p\)-value=0.435). In other words, based on the results of this ANOVA,
\(\mathcal{M}_2\) is preferred.

We can repeat the same process for \texttt{age} as well. We see there is
very clear evidence that a non-linear term is required for \texttt{age}.

    \begin{tcolorbox}[breakable, size=fbox, boxrule=1pt, pad at break*=1mm,colback=cellbackground, colframe=cellborder]
\prompt{In}{incolor}{38}{\boxspacing}
\begin{Verbatim}[commandchars=\\\{\}]
\PY{n}{gam\PYZus{}0} \PY{o}{=} \PY{n}{LinearGAM}\PY{p}{(}\PY{n}{year\PYZus{}term} \PY{o}{+}
                  \PY{n}{f\PYZus{}gam}\PY{p}{(}\PY{l+m+mi}{2}\PY{p}{,} \PY{n}{lam}\PY{o}{=}\PY{l+m+mi}{0}\PY{p}{)}\PY{p}{)}
\PY{n}{gam\PYZus{}linear} \PY{o}{=} \PY{n}{LinearGAM}\PY{p}{(}\PY{n}{l\PYZus{}gam}\PY{p}{(}\PY{l+m+mi}{0}\PY{p}{,} \PY{n}{lam}\PY{o}{=}\PY{l+m+mi}{0}\PY{p}{)} \PY{o}{+}
                       \PY{n}{year\PYZus{}term} \PY{o}{+}
                       \PY{n}{f\PYZus{}gam}\PY{p}{(}\PY{l+m+mi}{2}\PY{p}{,} \PY{n}{lam}\PY{o}{=}\PY{l+m+mi}{0}\PY{p}{)}\PY{p}{)}
\PY{n}{gam\PYZus{}0}\PY{o}{.}\PY{n}{fit}\PY{p}{(}\PY{n}{Xgam}\PY{p}{,} \PY{n}{y}\PY{p}{)}
\PY{n}{gam\PYZus{}linear}\PY{o}{.}\PY{n}{fit}\PY{p}{(}\PY{n}{Xgam}\PY{p}{,} \PY{n}{y}\PY{p}{)}
\PY{n}{anova\PYZus{}gam}\PY{p}{(}\PY{n}{gam\PYZus{}0}\PY{p}{,} \PY{n}{gam\PYZus{}linear}\PY{p}{,} \PY{n}{gam\PYZus{}full}\PY{p}{)}
\end{Verbatim}
\end{tcolorbox}

            \begin{tcolorbox}[breakable, size=fbox, boxrule=.5pt, pad at break*=1mm, opacityfill=0]
\prompt{Out}{outcolor}{38}{\boxspacing}
\begin{Verbatim}[commandchars=\\\{\}]
       deviance           df  deviance\_diff   df\_diff           F  \textbackslash{}
0  3.975443e+06  2991.000589            NaN       NaN         NaN
1  3.850247e+06  2990.000704  125196.137317  0.999884  101.270106
2  3.693143e+06  2987.007254  157103.978302  2.993450   42.447812

         pvalue
0           NaN
1  1.681120e-07
2  5.669414e-07
\end{Verbatim}
\end{tcolorbox}
        
    There is a (verbose) \texttt{summary()} method for the GAM fit.

    \begin{tcolorbox}[breakable, size=fbox, boxrule=1pt, pad at break*=1mm,colback=cellbackground, colframe=cellborder]
\prompt{In}{incolor}{39}{\boxspacing}
\begin{Verbatim}[commandchars=\\\{\}]
\PY{n}{gam\PYZus{}full}\PY{o}{.}\PY{n}{summary}\PY{p}{(}\PY{p}{)}
\end{Verbatim}
\end{tcolorbox}

    \begin{Verbatim}[commandchars=\\\{\}]
LinearGAM
===============================================
==========================================================
Distribution:                        NormalDist Effective DoF:
12.9927
Link Function:                     IdentityLink Log Likelihood:
-24117.907
Number of Samples:                         3000 AIC:
48263.7995
                                                AICc:
48263.94
                                                GCV:
1246.1129
                                                Scale:
1236.4024
                                                Pseudo R-Squared:
0.2928
================================================================================
==========================
Feature Function                  Lambda               Rank         EDoF
P > x        Sig. Code
================================= ==================== ============ ============
============ ============
s(0)                              [465.0491]           20           5.1
1.11e-16     ***
s(1)                              [2.1564]             7            4.0
8.10e-03     **
f(2)                              [0]                  5            4.0
1.11e-16     ***
intercept                                              1            0.0
1.11e-16     ***
================================================================================
==========================
Significance codes:  0 '***' 0.001 '**' 0.01 '*' 0.05 '.' 0.1 ' ' 1

WARNING: Fitting splines and a linear function to a feature introduces a model
identifiability problem
         which can cause p-values to appear significant when they are not.

WARNING: p-values calculated in this manner behave correctly for un-penalized
models or models with
         known smoothing parameters, but when smoothing parameters have been
estimated, the p-values
         are typically lower than they should be, meaning that the tests reject
the null too readily.
    \end{Verbatim}

    \begin{Verbatim}[commandchars=\\\{\}]
/var/folders/16/8y65\_zv174qgdp4ktlmpv12h0000gq/T/ipykernel\_70598/2135516388.py:1
: UserWarning: KNOWN BUG: p-values computed in this summary are likely much
smaller than they should be.

Please do not make inferences based on these values!

Collaborate on a solution, and stay up to date at:
github.com/dswah/pyGAM/issues/163

  gam\_full.summary()
    \end{Verbatim}

    We can make predictions from \texttt{gam} objects, just like from
\texttt{lm} objects, using the \texttt{predict()} method for the class
\texttt{gam}. Here we make predictions on the training set.

    \begin{tcolorbox}[breakable, size=fbox, boxrule=1pt, pad at break*=1mm,colback=cellbackground, colframe=cellborder]
\prompt{In}{incolor}{40}{\boxspacing}
\begin{Verbatim}[commandchars=\\\{\}]
\PY{n}{Yhat} \PY{o}{=} \PY{n}{gam\PYZus{}full}\PY{o}{.}\PY{n}{predict}\PY{p}{(}\PY{n}{Xgam}\PY{p}{)}
\end{Verbatim}
\end{tcolorbox}

    In order to fit a logistic regression GAM, we use \texttt{LogisticGAM()}
from \texttt{pygam}.

    \begin{tcolorbox}[breakable, size=fbox, boxrule=1pt, pad at break*=1mm,colback=cellbackground, colframe=cellborder]
\prompt{In}{incolor}{41}{\boxspacing}
\begin{Verbatim}[commandchars=\\\{\}]
\PY{n}{gam\PYZus{}logit} \PY{o}{=} \PY{n}{LogisticGAM}\PY{p}{(}\PY{n}{age\PYZus{}term} \PY{o}{+} 
                        \PY{n}{l\PYZus{}gam}\PY{p}{(}\PY{l+m+mi}{1}\PY{p}{,} \PY{n}{lam}\PY{o}{=}\PY{l+m+mi}{0}\PY{p}{)} \PY{o}{+}
                        \PY{n}{f\PYZus{}gam}\PY{p}{(}\PY{l+m+mi}{2}\PY{p}{,} \PY{n}{lam}\PY{o}{=}\PY{l+m+mi}{0}\PY{p}{)}\PY{p}{)}
\PY{n}{gam\PYZus{}logit}\PY{o}{.}\PY{n}{fit}\PY{p}{(}\PY{n}{Xgam}\PY{p}{,} \PY{n}{high\PYZus{}earn}\PY{p}{)}
\end{Verbatim}
\end{tcolorbox}

            \begin{tcolorbox}[breakable, size=fbox, boxrule=.5pt, pad at break*=1mm, opacityfill=0]
\prompt{Out}{outcolor}{41}{\boxspacing}
\begin{Verbatim}[commandchars=\\\{\}]
LogisticGAM(callbacks=[Deviance(), Diffs(), Accuracy()],
   fit\_intercept=True, max\_iter=100,
   terms=s(0) + l(1) + f(2) + intercept, tol=0.0001, verbose=False)
\end{Verbatim}
\end{tcolorbox}
        
    \begin{tcolorbox}[breakable, size=fbox, boxrule=1pt, pad at break*=1mm,colback=cellbackground, colframe=cellborder]
\prompt{In}{incolor}{42}{\boxspacing}
\begin{Verbatim}[commandchars=\\\{\}]
\PY{n}{fig}\PY{p}{,} \PY{n}{ax} \PY{o}{=} \PY{n}{subplots}\PY{p}{(}\PY{n}{figsize}\PY{o}{=}\PY{p}{(}\PY{l+m+mi}{8}\PY{p}{,} \PY{l+m+mi}{8}\PY{p}{)}\PY{p}{)}
\PY{n}{ax} \PY{o}{=} \PY{n}{plot\PYZus{}gam}\PY{p}{(}\PY{n}{gam\PYZus{}logit}\PY{p}{,} \PY{l+m+mi}{2}\PY{p}{)}
\PY{n}{ax}\PY{o}{.}\PY{n}{set\PYZus{}xlabel}\PY{p}{(}\PY{l+s+s1}{\PYZsq{}}\PY{l+s+s1}{Education}\PY{l+s+s1}{\PYZsq{}}\PY{p}{)}
\PY{n}{ax}\PY{o}{.}\PY{n}{set\PYZus{}ylabel}\PY{p}{(}\PY{l+s+s1}{\PYZsq{}}\PY{l+s+s1}{Effect on wage}\PY{l+s+s1}{\PYZsq{}}\PY{p}{)}
\PY{n}{ax}\PY{o}{.}\PY{n}{set\PYZus{}title}\PY{p}{(}\PY{l+s+s1}{\PYZsq{}}\PY{l+s+s1}{Partial dependence of wage on education}\PY{l+s+s1}{\PYZsq{}}\PY{p}{,}
             \PY{n}{fontsize}\PY{o}{=}\PY{l+m+mi}{20}\PY{p}{)}\PY{p}{;}
\PY{n}{ax}\PY{o}{.}\PY{n}{set\PYZus{}xticklabels}\PY{p}{(}\PY{n}{Wage}\PY{p}{[}\PY{l+s+s1}{\PYZsq{}}\PY{l+s+s1}{education}\PY{l+s+s1}{\PYZsq{}}\PY{p}{]}\PY{o}{.}\PY{n}{cat}\PY{o}{.}\PY{n}{categories}\PY{p}{,} \PY{n}{fontsize}\PY{o}{=}\PY{l+m+mi}{8}\PY{p}{)}\PY{p}{;}
\end{Verbatim}
\end{tcolorbox}

    \begin{center}
    \adjustimage{max size={0.9\linewidth}{0.9\paperheight}}{artifacts/latex/Ch07-nonlin-lab_files/artifacts/latex/Ch07-nonlin-lab_87_0.png}
    \end{center}
    { \hspace*{\fill} \\}
    
    The model seems to be very flat, with especially high error bars for the
first category. Let's look at the data a bit more closely.

    \begin{tcolorbox}[breakable, size=fbox, boxrule=1pt, pad at break*=1mm,colback=cellbackground, colframe=cellborder]
\prompt{In}{incolor}{43}{\boxspacing}
\begin{Verbatim}[commandchars=\\\{\}]
\PY{n}{pd}\PY{o}{.}\PY{n}{crosstab}\PY{p}{(}\PY{n}{Wage}\PY{p}{[}\PY{l+s+s1}{\PYZsq{}}\PY{l+s+s1}{high\PYZus{}earn}\PY{l+s+s1}{\PYZsq{}}\PY{p}{]}\PY{p}{,} \PY{n}{Wage}\PY{p}{[}\PY{l+s+s1}{\PYZsq{}}\PY{l+s+s1}{education}\PY{l+s+s1}{\PYZsq{}}\PY{p}{]}\PY{p}{)}
\end{Verbatim}
\end{tcolorbox}

            \begin{tcolorbox}[breakable, size=fbox, boxrule=.5pt, pad at break*=1mm, opacityfill=0]
\prompt{Out}{outcolor}{43}{\boxspacing}
\begin{Verbatim}[commandchars=\\\{\}]
education  1. < HS Grad  2. HS Grad  3. Some College  4. College Grad  \textbackslash{}
high\_earn
False               268         966              643              663
True                  0           5                7               22

education  5. Advanced Degree
high\_earn
False                     381
True                       45
\end{Verbatim}
\end{tcolorbox}
        
    We see that there are no high earners in the first category of
education, meaning that the model will have a hard time fitting. We will
fit a logistic regression GAM excluding all observations falling into
this category. This provides more sensible results.

To do so, we could subset the model matrix, though this will not remove
the column from \texttt{Xgam}. While we can deduce which column
corresponds to this feature, for reproducibility's sake we reform the
model matrix on this smaller subset.

    \begin{tcolorbox}[breakable, size=fbox, boxrule=1pt, pad at break*=1mm,colback=cellbackground, colframe=cellborder]
\prompt{In}{incolor}{44}{\boxspacing}
\begin{Verbatim}[commandchars=\\\{\}]
\PY{n}{only\PYZus{}hs} \PY{o}{=} \PY{n}{Wage}\PY{p}{[}\PY{l+s+s1}{\PYZsq{}}\PY{l+s+s1}{education}\PY{l+s+s1}{\PYZsq{}}\PY{p}{]} \PY{o}{==} \PY{l+s+s1}{\PYZsq{}}\PY{l+s+s1}{1. \PYZlt{} HS Grad}\PY{l+s+s1}{\PYZsq{}}
\PY{n}{Wage\PYZus{}} \PY{o}{=} \PY{n}{Wage}\PY{o}{.}\PY{n}{loc}\PY{p}{[}\PY{o}{\PYZti{}}\PY{n}{only\PYZus{}hs}\PY{p}{]}
\PY{n}{Xgam\PYZus{}} \PY{o}{=} \PY{n}{np}\PY{o}{.}\PY{n}{column\PYZus{}stack}\PY{p}{(}\PY{p}{[}\PY{n}{Wage\PYZus{}}\PY{p}{[}\PY{l+s+s1}{\PYZsq{}}\PY{l+s+s1}{age}\PY{l+s+s1}{\PYZsq{}}\PY{p}{]}\PY{p}{,}
                         \PY{n}{Wage\PYZus{}}\PY{p}{[}\PY{l+s+s1}{\PYZsq{}}\PY{l+s+s1}{year}\PY{l+s+s1}{\PYZsq{}}\PY{p}{]}\PY{p}{,}
                         \PY{n}{Wage\PYZus{}}\PY{p}{[}\PY{l+s+s1}{\PYZsq{}}\PY{l+s+s1}{education}\PY{l+s+s1}{\PYZsq{}}\PY{p}{]}\PY{o}{.}\PY{n}{cat}\PY{o}{.}\PY{n}{codes}\PY{o}{\PYZhy{}}\PY{l+m+mi}{1}\PY{p}{]}\PY{p}{)}
\PY{n}{high\PYZus{}earn\PYZus{}} \PY{o}{=} \PY{n}{Wage\PYZus{}}\PY{p}{[}\PY{l+s+s1}{\PYZsq{}}\PY{l+s+s1}{high\PYZus{}earn}\PY{l+s+s1}{\PYZsq{}}\PY{p}{]}
\end{Verbatim}
\end{tcolorbox}

    In the second-to-last line above, we subtract one from the codes of the
category, due to a bug in \texttt{pygam}. It just relabels the education
values and hence has no effect on the fit.

We now fit the model.

    \begin{tcolorbox}[breakable, size=fbox, boxrule=1pt, pad at break*=1mm,colback=cellbackground, colframe=cellborder]
\prompt{In}{incolor}{45}{\boxspacing}
\begin{Verbatim}[commandchars=\\\{\}]
\PY{n}{gam\PYZus{}logit\PYZus{}} \PY{o}{=} \PY{n}{LogisticGAM}\PY{p}{(}\PY{n}{age\PYZus{}term} \PY{o}{+}
                         \PY{n}{year\PYZus{}term} \PY{o}{+}
                         \PY{n}{f\PYZus{}gam}\PY{p}{(}\PY{l+m+mi}{2}\PY{p}{,} \PY{n}{lam}\PY{o}{=}\PY{l+m+mi}{0}\PY{p}{)}\PY{p}{)}
\PY{n}{gam\PYZus{}logit\PYZus{}}\PY{o}{.}\PY{n}{fit}\PY{p}{(}\PY{n}{Xgam\PYZus{}}\PY{p}{,} \PY{n}{high\PYZus{}earn\PYZus{}}\PY{p}{)}
\end{Verbatim}
\end{tcolorbox}

            \begin{tcolorbox}[breakable, size=fbox, boxrule=.5pt, pad at break*=1mm, opacityfill=0]
\prompt{Out}{outcolor}{45}{\boxspacing}
\begin{Verbatim}[commandchars=\\\{\}]
LogisticGAM(callbacks=[Deviance(), Diffs(), Accuracy()],
   fit\_intercept=True, max\_iter=100,
   terms=s(0) + s(1) + f(2) + intercept, tol=0.0001, verbose=False)
\end{Verbatim}
\end{tcolorbox}
        
    Let's look at the effect of \texttt{education}, \texttt{year} and
\texttt{age} on high earner status now that we've removed those
observations.

    \begin{tcolorbox}[breakable, size=fbox, boxrule=1pt, pad at break*=1mm,colback=cellbackground, colframe=cellborder]
\prompt{In}{incolor}{46}{\boxspacing}
\begin{Verbatim}[commandchars=\\\{\}]
\PY{n}{fig}\PY{p}{,} \PY{n}{ax} \PY{o}{=} \PY{n}{subplots}\PY{p}{(}\PY{n}{figsize}\PY{o}{=}\PY{p}{(}\PY{l+m+mi}{8}\PY{p}{,} \PY{l+m+mi}{8}\PY{p}{)}\PY{p}{)}
\PY{n}{ax} \PY{o}{=} \PY{n}{plot\PYZus{}gam}\PY{p}{(}\PY{n}{gam\PYZus{}logit\PYZus{}}\PY{p}{,} \PY{l+m+mi}{2}\PY{p}{)}
\PY{n}{ax}\PY{o}{.}\PY{n}{set\PYZus{}xlabel}\PY{p}{(}\PY{l+s+s1}{\PYZsq{}}\PY{l+s+s1}{Education}\PY{l+s+s1}{\PYZsq{}}\PY{p}{)}
\PY{n}{ax}\PY{o}{.}\PY{n}{set\PYZus{}ylabel}\PY{p}{(}\PY{l+s+s1}{\PYZsq{}}\PY{l+s+s1}{Effect on wage}\PY{l+s+s1}{\PYZsq{}}\PY{p}{)}
\PY{n}{ax}\PY{o}{.}\PY{n}{set\PYZus{}title}\PY{p}{(}\PY{l+s+s1}{\PYZsq{}}\PY{l+s+s1}{Partial dependence of high earner status on education}\PY{l+s+s1}{\PYZsq{}}\PY{p}{,} \PY{n}{fontsize}\PY{o}{=}\PY{l+m+mi}{20}\PY{p}{)}\PY{p}{;}
\PY{n}{ax}\PY{o}{.}\PY{n}{set\PYZus{}xticklabels}\PY{p}{(}\PY{n}{Wage}\PY{p}{[}\PY{l+s+s1}{\PYZsq{}}\PY{l+s+s1}{education}\PY{l+s+s1}{\PYZsq{}}\PY{p}{]}\PY{o}{.}\PY{n}{cat}\PY{o}{.}\PY{n}{categories}\PY{p}{[}\PY{l+m+mi}{1}\PY{p}{:}\PY{p}{]}\PY{p}{,}
                   \PY{n}{fontsize}\PY{o}{=}\PY{l+m+mi}{8}\PY{p}{)}\PY{p}{;}
\end{Verbatim}
\end{tcolorbox}

    \begin{center}
    \adjustimage{max size={0.9\linewidth}{0.9\paperheight}}{artifacts/latex/Ch07-nonlin-lab_files/artifacts/latex/Ch07-nonlin-lab_95_0.png}
    \end{center}
    { \hspace*{\fill} \\}
    
    \begin{tcolorbox}[breakable, size=fbox, boxrule=1pt, pad at break*=1mm,colback=cellbackground, colframe=cellborder]
\prompt{In}{incolor}{47}{\boxspacing}
\begin{Verbatim}[commandchars=\\\{\}]
\PY{n}{fig}\PY{p}{,} \PY{n}{ax} \PY{o}{=} \PY{n}{subplots}\PY{p}{(}\PY{n}{figsize}\PY{o}{=}\PY{p}{(}\PY{l+m+mi}{8}\PY{p}{,} \PY{l+m+mi}{8}\PY{p}{)}\PY{p}{)}
\PY{n}{ax} \PY{o}{=} \PY{n}{plot\PYZus{}gam}\PY{p}{(}\PY{n}{gam\PYZus{}logit\PYZus{}}\PY{p}{,} \PY{l+m+mi}{1}\PY{p}{)}
\PY{n}{ax}\PY{o}{.}\PY{n}{set\PYZus{}xlabel}\PY{p}{(}\PY{l+s+s1}{\PYZsq{}}\PY{l+s+s1}{Year}\PY{l+s+s1}{\PYZsq{}}\PY{p}{)}
\PY{n}{ax}\PY{o}{.}\PY{n}{set\PYZus{}ylabel}\PY{p}{(}\PY{l+s+s1}{\PYZsq{}}\PY{l+s+s1}{Effect on wage}\PY{l+s+s1}{\PYZsq{}}\PY{p}{)}
\PY{n}{ax}\PY{o}{.}\PY{n}{set\PYZus{}title}\PY{p}{(}\PY{l+s+s1}{\PYZsq{}}\PY{l+s+s1}{Partial dependence of high earner status on year}\PY{l+s+s1}{\PYZsq{}}\PY{p}{,}
             \PY{n}{fontsize}\PY{o}{=}\PY{l+m+mi}{20}\PY{p}{)}\PY{p}{;}
\end{Verbatim}
\end{tcolorbox}

    \begin{center}
    \adjustimage{max size={0.9\linewidth}{0.9\paperheight}}{artifacts/latex/Ch07-nonlin-lab_files/artifacts/latex/Ch07-nonlin-lab_96_0.png}
    \end{center}
    { \hspace*{\fill} \\}
    
    \begin{tcolorbox}[breakable, size=fbox, boxrule=1pt, pad at break*=1mm,colback=cellbackground, colframe=cellborder]
\prompt{In}{incolor}{48}{\boxspacing}
\begin{Verbatim}[commandchars=\\\{\}]
\PY{n}{fig}\PY{p}{,} \PY{n}{ax} \PY{o}{=} \PY{n}{subplots}\PY{p}{(}\PY{n}{figsize}\PY{o}{=}\PY{p}{(}\PY{l+m+mi}{8}\PY{p}{,} \PY{l+m+mi}{8}\PY{p}{)}\PY{p}{)}
\PY{n}{ax} \PY{o}{=} \PY{n}{plot\PYZus{}gam}\PY{p}{(}\PY{n}{gam\PYZus{}logit\PYZus{}}\PY{p}{,} \PY{l+m+mi}{0}\PY{p}{)}
\PY{n}{ax}\PY{o}{.}\PY{n}{set\PYZus{}xlabel}\PY{p}{(}\PY{l+s+s1}{\PYZsq{}}\PY{l+s+s1}{Age}\PY{l+s+s1}{\PYZsq{}}\PY{p}{)}
\PY{n}{ax}\PY{o}{.}\PY{n}{set\PYZus{}ylabel}\PY{p}{(}\PY{l+s+s1}{\PYZsq{}}\PY{l+s+s1}{Effect on wage}\PY{l+s+s1}{\PYZsq{}}\PY{p}{)}
\PY{n}{ax}\PY{o}{.}\PY{n}{set\PYZus{}title}\PY{p}{(}\PY{l+s+s1}{\PYZsq{}}\PY{l+s+s1}{Partial dependence of high earner status on age}\PY{l+s+s1}{\PYZsq{}}\PY{p}{,} \PY{n}{fontsize}\PY{o}{=}\PY{l+m+mi}{20}\PY{p}{)}\PY{p}{;}
\end{Verbatim}
\end{tcolorbox}

    \begin{center}
    \adjustimage{max size={0.9\linewidth}{0.9\paperheight}}{artifacts/latex/Ch07-nonlin-lab_files/artifacts/latex/Ch07-nonlin-lab_97_0.png}
    \end{center}
    { \hspace*{\fill} \\}
    
    \hypertarget{local-regression}{%
\subsection{Local Regression}\label{local-regression}}

We illustrate the use of local regression using the \texttt{lowess()}
function from \texttt{sm.nonparametric}. Some implementations of GAMs
allow terms to be local regression operators; this is not the case in
\texttt{pygam}.

Here we fit local linear regression models using spans of 0.2 and 0.5;
that is, each neighborhood consists of 20\% or 50\% of the observations.
As expected, using a span of 0.5 is smoother than 0.2.

    \begin{tcolorbox}[breakable, size=fbox, boxrule=1pt, pad at break*=1mm,colback=cellbackground, colframe=cellborder]
\prompt{In}{incolor}{49}{\boxspacing}
\begin{Verbatim}[commandchars=\\\{\}]
\PY{n}{lowess} \PY{o}{=} \PY{n}{sm}\PY{o}{.}\PY{n}{nonparametric}\PY{o}{.}\PY{n}{lowess}
\PY{n}{fig}\PY{p}{,} \PY{n}{ax} \PY{o}{=} \PY{n}{subplots}\PY{p}{(}\PY{n}{figsize}\PY{o}{=}\PY{p}{(}\PY{l+m+mi}{8}\PY{p}{,}\PY{l+m+mi}{8}\PY{p}{)}\PY{p}{)}
\PY{n}{ax}\PY{o}{.}\PY{n}{scatter}\PY{p}{(}\PY{n}{age}\PY{p}{,} \PY{n}{y}\PY{p}{,} \PY{n}{facecolor}\PY{o}{=}\PY{l+s+s1}{\PYZsq{}}\PY{l+s+s1}{gray}\PY{l+s+s1}{\PYZsq{}}\PY{p}{,} \PY{n}{alpha}\PY{o}{=}\PY{l+m+mf}{0.5}\PY{p}{)}
\PY{k}{for} \PY{n}{span} \PY{o+ow}{in} \PY{p}{[}\PY{l+m+mf}{0.2}\PY{p}{,} \PY{l+m+mf}{0.5}\PY{p}{]}\PY{p}{:}
    \PY{n}{fitted} \PY{o}{=} \PY{n}{lowess}\PY{p}{(}\PY{n}{y}\PY{p}{,}
                    \PY{n}{age}\PY{p}{,}
                    \PY{n}{frac}\PY{o}{=}\PY{n}{span}\PY{p}{,}
                    \PY{n}{xvals}\PY{o}{=}\PY{n}{age\PYZus{}grid}\PY{p}{)}
    \PY{n}{ax}\PY{o}{.}\PY{n}{plot}\PY{p}{(}\PY{n}{age\PYZus{}grid}\PY{p}{,}
            \PY{n}{fitted}\PY{p}{,}
            \PY{n}{label}\PY{o}{=}\PY{l+s+s1}{\PYZsq{}}\PY{l+s+si}{\PYZob{}:.1f\PYZcb{}}\PY{l+s+s1}{\PYZsq{}}\PY{o}{.}\PY{n}{format}\PY{p}{(}\PY{n}{span}\PY{p}{)}\PY{p}{,}
            \PY{n}{linewidth}\PY{o}{=}\PY{l+m+mi}{4}\PY{p}{)}
\PY{n}{ax}\PY{o}{.}\PY{n}{set\PYZus{}xlabel}\PY{p}{(}\PY{l+s+s1}{\PYZsq{}}\PY{l+s+s1}{Age}\PY{l+s+s1}{\PYZsq{}}\PY{p}{,} \PY{n}{fontsize}\PY{o}{=}\PY{l+m+mi}{20}\PY{p}{)}
\PY{n}{ax}\PY{o}{.}\PY{n}{set\PYZus{}ylabel}\PY{p}{(}\PY{l+s+s1}{\PYZsq{}}\PY{l+s+s1}{Wage}\PY{l+s+s1}{\PYZsq{}}\PY{p}{,} \PY{n}{fontsize}\PY{o}{=}\PY{l+m+mi}{20}\PY{p}{)}\PY{p}{;}
\PY{n}{ax}\PY{o}{.}\PY{n}{legend}\PY{p}{(}\PY{n}{title}\PY{o}{=}\PY{l+s+s1}{\PYZsq{}}\PY{l+s+s1}{span}\PY{l+s+s1}{\PYZsq{}}\PY{p}{,} \PY{n}{fontsize}\PY{o}{=}\PY{l+m+mi}{15}\PY{p}{)}\PY{p}{;}
\end{Verbatim}
\end{tcolorbox}

    \begin{center}
    \adjustimage{max size={0.9\linewidth}{0.9\paperheight}}{artifacts/latex/Ch07-nonlin-lab_files/artifacts/latex/Ch07-nonlin-lab_99_0.png}
    \end{center}
    { \hspace*{\fill} \\}
    
    


    % Add a bibliography block to the postdoc
    
    
    
\end{document}
